%  LaTeX support: latex@mdpi.com 
%  For support, please attach all files needed for compiling as well as the log file, and specify your operating system, LaTeX version, and LaTeX editor.

%=================================================================
\documentclass[land,article,submit,pdftex,moreauthors]{Definitions/mdpi} 

%=================================================================
%----------
% journal
%----------
% Choose between the following MDPI journals:

% MDPI internal commands
\firstpage{1} 
\makeatletter 
\setcounter{page}{\@firstpage} 
\makeatother
\pubvolume{1}
\issuenum{1}
\articlenumber{0}
\pubyear{2023}
\copyrightyear{2023}
%\externaleditor{Academic Editor: Firstname Lastname}
\datereceived{22 September 2023} 
\daterevised{} 
\dateaccepted{} 
\datepublished{} 
\hreflink{https://doi.org/} % If needed use \linebreak
%\doinum{}

%=================================================================
% Add packages and commands here. The following packages are loaded in our class file: fontenc, inputenc, calc, indentfirst, fancyhdr, graphicx, epstopdf, lastpage, ifthen, lineno, float, amsmath, setspace, enumitem, mathpazo, booktabs, titlesec, etoolbox, tabto, xcolor, soul, multirow, microtype, tikz, totcount, changepage, attrib, upgreek, cleveref, amsthm, hyphenat, natbib, hyperref, footmisc, url, geometry, newfloat, caption

\graphicspath{{figures/}} % path to folder with figures
\usepackage{subcaption}
\usepackage{listings}
\usepackage{xcolor}
\usepackage{gensymb} % for \textdegree
\usepackage{tabularx}
\usepackage{makecell}
\usepackage{multirow}
\usepackage{verbatim}

%=================================================================
% Full title of the paper (Capitalized)
\Title{Monitoring Seasonal Fluctuations in Sabkhas of Tunisia by Multispectral Image Analysis Using GRASS GIS}

% MDPI internal command: Title for citation in the left column
\TitleCitation{Monitoring Seasonal Fluctuations in Sabkhas of Tunisia by Multispectral Image Analysis Using GRASS GIS}

% Author Orchid ID: enter ID or remove command
\newcommand{\orcidauthorA}{0000-0002-5759-1089} % Polina Add \orcidA{} behind the author's name

% Authors, for the paper (add full first names)
\Author{Polina Lemenkova $^{1}$*\orcidA{}}

%\longauthorlist{yes}

% MDPI internal command: Authors, for metadata in PDF
\AuthorNames{Polina Lemenkova}

% MDPI internal command: Authors, for citation in the left column
\AuthorCitation{Lemenkova, P.}
% If this is a Chicago style journal: Lastname, Firstname, Firstname Lastname, and Firstname Lastname.

% Affiliations / Addresses (Add [1] after \address if there is only one affiliation.)
\address{%
$^{1}$ \quad Department of Geoinformatics, Universität Salzburg, Schillerstraße 30, A-5020 Salzburg, Austria}

% Contact information of the corresponding author
\corres{Correspondence: polina.lemenkova@plus.ac.at; Tel.: +43-677-6173-2772 (P.L.)}

% Current address and/or shared authorship
%\firstnote{Current address: Affiliation 3.} 

% Abstract (Do not insert blank lines, i.e. \\) 
\abstract{This study documents the variations of the saline lakes (a.k.a. sebkhet or sabkha) in Tunisia due to the climate effects. Located in arid environment of north Africa, salt lakes experience changes in volume and extent due to extreme temperatures: during summer heat period, lakes dry out almost entirely and become a crust of salt and minerals due to the evaporated moisture. Lakes in Tunisia are important components of the environment playing a crucial role as biodiversity maintenance, habitats for rare species in nearby oases, rare sources of water, flood and fishery, regulators of temperature, hydrology, and carbon fixation. Remote sensing data are valuable data source to monitoring lakes since variations in the extent of lakes are visible from space. In this study, changes in salinisation and water extent of the salt pans of Tunisia were evaluated using a series of 12 Landsat 8-9 OLI/TIRS images processed with GRASS GIS software. The study area included four salt lakes of north Tunisia in the two regions of the Gulf of Hammamet and Gulf of Gabès: 1) Sebkhet de Sidi el Hani (Sousse Governorate), 2) Sebkha de Moknine (Mahdia Governorate), 3) Sebkhet El Rharra and 4) Sebkhet en Noual (Sfax). The scenes were analysed for a periods 2017--2023 on months February, April and July for each year. Spatio-temporal changes and their climate-environmental driving forces were analysed. The results were interpreted and the highest changes detected by accuracy assessment, computing the class separability matrices, evaluating the means and standard deviation for each band, and plotting the reject probability maps. Multi-temporal changes in land cover classes are reported for each image. The results demonstrated changes in salt lakes were determined for winter/spring/summer months as detected changes in water/land/salt/sand/vegetation areas. This paper contributes to the environmental monitoring of Tunisian landscapes and analysis of climate effects on landscapes of north Africa.}

% Keywords
\keyword{Remote sensing; image processing; Africa; Sahara; salt lakes; desert environment; arid climate; Tunisia; image analysis; Landsat} 

% The fields PACS, MSC, and JEL may be left empty or commented out if not applicable
\PACS{91.10.Da; 91.10.Jf; 91.10.Sp; 91.10.Xa; 96.25.Vt; 91.10.Fc; 95.40.+s; 95.75.Qr; 95.75.Rs; 42.68.Wt}
\MSC{86A30; 86-08; 86A99; 86A04}
\JEL{Y91; Q20; Q24; Q23; Q3; Q01; R11; O44; O13; Q5; Q51; Q55; N57; C6; C61}

\begin{document}

%----------- section ----------->
\section{Introduction}

\subsection{Background}

Remote sensing (RS) data analysis is a fundamental issue in environmental studies and ecosystem monitoring. Processing RS data is especially essential for multi-temporal analysis of land cover changes \cite{BOUSSETTA2022102564} and application of environmental geomorphology, which are subject to seasonal changes and shrinkage due to climate effects \cite{JonesMillington}. Climate effects can cause desertification \cite{Pontanier,info14040249} and threaten balance of biodiversity \cite{YAICHEACHOUR2023501}. Therefore, monitoring environmental properties of Saharan landscapes by RS data plays an essential role in analysis and decision making on environmental sustainability, which results in many approaches to the effective methods of image processing for environmental tasks \cite{Monget,Lemenkova202227,CAMPOS2018211}. 

The key factors in effective approaches to satellite image processing include the type of data and the functionality of software \cite{rs14236017}. Another important factor consists in the quality of image, coverage extent, cloudiness and haze. Major image types widely used for effective environmental monitoring and include Landsat \cite{BOUSSEMA2023185,SMIDA2022104643,jmse11040871}, Advanced Very High Resolution Radiometer (AVHRR) \cite{Bryant}, Google Earth \cite{LI2022252}, Moderate Resolution Imaging Spectroradiometer (MODIS) \cite{TOUHAMI2022103804,RHIF2022101596}, and the combinations thereof \cite{REBILLARD1984133}. Examples of very high resolution data include such images as Quickbird \cite{CHOUARI2021443}, Satellite Pour l’Observation de la Terre (SPOT) and IKONOS \cite{VELA2008591}. Certainly, although diverse images can be used in various tasks as data sources for environmental monitoring and mapping, the general approach of image analysis relies on using imagery as Landsat or other existing sensor types to extract information through image processing \cite{LiangWang}. 

Selecting a suitable GIS is important for environmental mapping of Chad Lake, since its functionality is necessary for data analysis and cartographic techniques. A comprehensive survey is available on the existing free GIS \cite{LEIDIG201549,CHEN2010253}, recent trends in the development of cartography \cite{WangWu} and new directions and perspectives in GIS \cite{Lietal2022}, while this section covers the most relevant GIS applications for environmental mapping. The efficient solutions for satellite image processing are presented by Earth Resources Data Analysis System (ERDAS) Imagine with applications of environmental monitoring \cite{FALAKI2020363} with many existing cases due to the functionality and availability of diverse modules aimed at raster data processing. The Global Mapper GIS, developed by the United States Geological Survey (USGS), includes a variety of tools including the Light Detection and Ranging (LiDAR) handling, 3D rendering, watershed delineation to process geological and environmental data. Probably the most-well known and widely distributed commercial GIS is the ArcGIS which process well both vector and raster data \cite{AlDjazouli}, supports satellite images processing \cite{SESNIE20082145}, enables hydrological catchment and basin analysis \cite{Nkiaka}, terrain DEM modelling for geospatial data analysis \cite{Kay} and spatial data visualization. 

Nevertheless, in order to reduce the uncertainty and increase the effectiveness of environmental monitoring, a lot of approaches to RS data are proposed as an alternative to processing climate data, some of them using Landsat products \cite{Lemenkova202315} and their integration with Google Earth Engine platform \cite{rs15133257}, or mixed approaches with GIS-based mapping \cite{SAIDI20091683,ANANE201236,Lemenkova202217}, and some of them using hydrological datasets for evaluating physicochemical parameters of groundwater in inflow \cite{BENBRAHIM2021104224}. Other examples include evaluating the content of soil organic carbon stocks in topsoils using bioclimatic data \cite{BAHRI2022e00561} or monitoring agricultural and irrigated lands and their behaviour during drought periods \cite{MOHAMED2023e01766}. 

Alternative approaches to environmental monitoring may be used as ways of optimised data modelling in various ecological tasks, to highlight correlations between climate and environmental responses in arid landscape \cite{BESSER2021104134} and to evaluate regional hydrogeological processes for a specific soil type in Sahara desert which affects the geochemical composition of groundwater that influences surface water quality and vegetation patterns in Tunisia \cite{HAMDI2021102974,RAFAATRIGUI2021104121}. At the same time, various landscapes have different mechanisms of ecosystem functionality controlling the rates of vegetation response to droughts and reaction in soil materials and phases of dissolution–precipitation during salinisation \cite{NAVARROTORRE2023105397}. Therefore, understanding the parameters of these processes, the effects of seasonal rise in temperature and soil performance as a sediment background in saline lakes of Tunisia during heat periods is of great importance for environmental monitoring of Tunisian sector of Sahara.

\subsection{Study Area}

This study focuses on the salt lakes located in the north Sahara, Tunisia, Figure \ref{fig01}. Salt lakes, or sebkhas, of Tunisia present a series of saline depressions that are characterised by variability of water: during the winter, sebkhas receive rare water inflow from occasional rains while during summer months the lakes dry out almost completely and become the saline crystallised surfaces \cite{berthon1970saumures,m2001valorisation,Gautier}. The geographical location of these sebkhas is related to the environmental historical development of north Sahara with related aeolian and fluvial processes \cite{ALI20231,MWIRICHIA2023269}. Gradual deflation and geomorphic erosion lead to the formation of low individual basins and pans where current salt lakes are located \cite{Blackwelder140}. The sebhkas are mostly located in topographic depressions which remain as geomorphic landforms from the earlier palaeolacustrine basins, palaeodrainages, interdunes, and coastal plains \cite{GOUDIE19951}. The formation of these depressions is strongly connected to the tectonic movements in Quaternary which shaped present morphology of sebkhas \cite{Abbes,Louhaichi,Ballais}. 

The lagoon facies of the selected lakes evident the stratigraphic distribution of gypsum clays and laminated gypsum lithofacies from the Neogene period (Mio--Pliocene and Quaternary evaporites) and Cretaceous sediments \cite{BARBIERI2006726,SWEZEY1996335}. This supports the hypothesis of earlier existence of the Sahara Sea with endorheic basins that nowadays become a series of depressions of the saline lakes in north Africa \cite{Coque}. 

\begin{figure}[H]
	\centering
		\includegraphics[width=13.0 cm]{Fig_01a.jpg}
\caption{Two segments of the study area on a topographic map of Tunisia: the Gulf of Hammamet (northern part) and the Gulf of Gabès (southern part). Software: GMT. Map source: author.}\label{fig01}
\end{figure}

The effects of climate and seasonal meteorological fluctuations on soil and vegetation setting of the Saharan landscapes are explained by changes in temperature which triggers processes of evaporation, i.e., the transform of liquid water in lakes and ephemeral river flows into gaseous vapour \cite{HAMMI2003215}. Natural factors are augmented by the anthropogenic activities that increased the intensity of land cover changes \cite{RODRIGUEZCABALLERO2017146}. High evaporation during the summer months results in a total desiccation of these ephemeral lakes and lead to a formation of the thick layer of minerals and salt on the bottom of the lake basins \cite{BRYANT1994269}. 

Lakes in Tunisia are important components of the dryland environment of Sahara. They are rare sources of water, flood and fishery, they serve as regulators of temperature, hydrology, and they support carbon fixation for environmental sustainability \cite{NIGAM2020201}. Moreover, they play a crucial role for maintaining biodiversity \cite{Neji}, they serve as habitats for rare African species and birds in nearby oases \cite{Baduel,RADDADI2022104806,Kamel} and habitats for halophilic species \cite{NEIFAR20191802}. Such phenomena were revealed in studies on environmental monitoring which evaluated landscape changes in Tunisia. For instance, the reaction of landscapes and vegetation responses to droughts were modelled with regard to the temperature fluctuations and climatic effects in Sahara \cite{HANAFI2008557,PAUSATA2020235,Lemenkova202313}. At the same time, the increase of temperature and reduced precipitation observed during heat of summer months lead to lack of surface water which is the primary factor of desertification and aridification of landscapes in Sahara \cite{OSWALD2023369,DODORICO2013326}. 

\subsection{Related work}

In this study, we explore the advanced cartographic approaches to map changes in saline lakes of Tunisia using satellite images to better evaluate the effects from climate and seasonal fluctuations in temperature on salinisation of sebkhas in north Sahara. The success in satellite image processing can be evaluated through classified and mapped difference in land cover types showing environmental characteristics of the soil and land categories which indicate land cover changes  \cite{Shoshany}. The traditional methods of evaluation of land cover change include the fieldwork and topographic measurements and the overlay of the existing vector layers using GIS. However, such methods have some drawbacks in that, the organisation of the fieldwork measurements is required for the whole process of land monitoring, surveying and topographic mapping \cite{JENDOUBI2020100507}. This might be hard and costly to implement in the remote areas of Sahara desert. Besides, traditional geographic fieldwork measurements provide information of land categories as an discrete values measured by the technical surveyors without data on complex patterns of the landscapes. 

An alternative approach is presented by RS data measurements and processing which can be used as descriptors of land cover types and partial replacement of the traditional topographic fieldwork methods. Thus, the use of the RS integral assessment of land cover patterns are possible using satellite images with large extent of the covered area \cite{ELGHOUL2023106733}. Thus, satellite images provide direct information on land categories due to the relationship between spectral reflectance and the environmental properties of land categories visible and interpreted on the spaceborne images \cite{Allouche,Ruiz,Ray}. As a general rule, vegetation and water areas can be identified using a combination of Red/NIR channels since the reflectance in NIR increase for vegetation along with the absorbed signal in Red bands which helps in discrimination between water and vegetation areas and during image analysis. Such RS properties are widely used for calculation of vegetation indices. 

The remote sensing approach is based on measuring the spectral reflectance of the pixels on the satellite image which represent diverse land cover types of Earth visible on the spaceborne images. Spectral reflectance of the pixels in various wavebands of the multispectral images indicates spectral characteristics of different objects on the Earth. This fundamental characteristics enables to use images for detecting land cover types and image interpretation after compensation and correcting the image from atmospheric and air-water effects. Therefore, evaluating land cover properties by using satellite images and land cover changes can be performed using times series approach to problem solving, i.e., interpretation of several satellite images taken on the same target area in different seasons of year or a sequence of several years. 

Remote sensing methods emerge as a fundamental technique of Earth Observation in a wide variety of geospatial and environmental applications including case studies on landscape mosaics \cite{BIARNES2021103281}, monitoring sedimentation processes \cite{TOWNSHEND1989177}, dynamics of sandy dunes in the Sahara \cite{Afrasinei}, geomorphological modelling \cite{Desprats} and monitoring coastal regions \cite{AMROUNI201973,Jaquet}. Such generic approach of remote sensing has been noted earlier \cite{RAMPHERI2023103359} with mentioned advantages of the use of spaceborne data. These include the high sensitivity of satellite images which enables to detect small-size patches of landscapes, regularity of survey by various space sensors and flexibility of data coverage which can be finely adjusted for practically any study area using search criteria in image selection and threshold parameters (cloudiness, image source, data, etc). 

When applied to a series of images, processing RS data can effectively discover the environmental properties of  landscapes \cite{Toujani,Quarmby,Gelebo}. Moreover, RS data processing and satellite image analysis enables to quantify the landscape patches within the target landscape through the interpretation of classes indicating land categories and vegetation types \cite{Benedetti,Godinho} and to determine climate effects in lacustrine or coastal environments \cite{Souissi}. Finally, other advantages of the RS data processing is that it is a resource- and time-effective method indicating dynamics in landscapes through evaluation of changes with time visible on a time series of satellite images. 

\subsection{Objectives and Motivation}

Inspired by the existing examples of the RS data analysis in environmental monitoring \cite{f9020059,Lemenkova202229,Fathalli,Abbas,Henchiri}, this study integrated the advanced methods of the GRASS GIS with applied techniques of its image processing modules for assessment and mapping of changes in landscapes of saline lakes in northern Tunisia, Figure \ref{fig01}. The main purpose of the current paper is to provide an advanced RS and cartographic methods that can assess the nonlinear behaviour of salt lakes in response to seasonal fluctuations in temperature with time indicating the development of salinisation. Importantly, the study presented here is not intended to suggest that RS data replace the hydrological datasets. In contrast, these images can supplement the fieldwork survey in the sabkhas and environmental tests on the saline soil of Tunisia. 

In this paper, the integration of image processing workflow with cartographic and geologic maps  provides an effective mechanism by which the analysis of environmental changes with regard to inter-year and multi-decades meteorological variations can be performed. The goal of using the RS data is to qualitatively evaluate and map changes in land cover types caused by seasonal fluctuations in temperature during winter/spring/summer months and in the process of evaporation. In this way, satellite image processing is the effective means by which the time series of the multispectral images of the Landsat OLI/TIRS sensor can be analysed for landscape monitoring. Data obtained from the clustering of the satellite images was used for modelling land cover types using classification. 

The cartographic experiments on image processing were performed on two selected study areas in north Tunisia in the coastal regions of the Gulf of Gabès (salt lake Sebkhet en Noual) and the Gulf of Hammamet (salt lakes Sebkhet de Sidi el Hani, Sebkha de Moknine and Sebkhet El Rharra). Using a comparison of several images classified for various seasons and time periods, this study is evaluating the properties of the landscapes using maximum-likelihood discriminant analysis classifier of GRASS GIS. The advanced methods of evaluation of landscape properties are built on the growing demand for the time- and cost-effective approaches to environmental monitoring in changing climate of Tunisia, saving agricultural resources, and taking decisions of farmers and environmental planners to make use of the renewable resources in the salt lake surrounding considering regional and local climate and environmental settings of north Tunisia.

%----------- section ----------->
\section{Materials and Methods}

%----------- subsection ----------->
\subsection{Data}

The materials include the 12 Landsat 8-9 OLI/TIRS images collected from the USGS Earth Explorer repository.The USGS Identifier (ID) of the scenes of the satellite images are summarised in Table \ref{tab01}. Within the dataset, 6 images cover the region of the Gulf of Hammamet (Figure \ref{fig03}), and 6 images cover the area of the Gulf of Gabès (Figure \ref{fig04}). 

\begin{table}[H]
\caption{Identifiers (ID) of the 12 Landsat 8-9 OLI/TIRS images in the EarthExplorer USGS repository}.
\label{tab01}
\begin{adjustwidth}{-\extralength}{-2.5cm}
\newcolumntype{L}{>{\centering\scriptsize\arraybackslash}l}
\setlength\extrarowheight{-0.8cm}
\footnotesize
%\begin{tabularx}{\fulllength}{p{1.5cm}|p{6.5cm}|p{6.5cm}|p{1.9cm}}
\begin{tabularx}{\fulllength}{p{1.4cm}|lll}
			\toprule
			\textbf{Date} & \textbf{Landsat Product Identifier L1} & \textbf{Landsat Product Identifier L2} & \textbf{Scene ID} \\
			\midrule
			\multicolumn{4}{c}{\textbf{Gulf of Hammamet region}}\\
			\midrule
			2023/07/15 & \makecell{LC09\_L1TP\_191035\_20230715\_20230715\_02\_T1} & \makecell{LC09\_L2SP\_191035\_20230715\_20230717\_02\_T1} & LC91910352023196LGN00 \\	
			2023/04/10 & \makecell{LC09\_L1TP\_191035\_20230410\_20230410\_02\_T1} & \makecell{LC09\_L2SP\_191035\_20230410\_20230412\_02\_T1} & LC91910352023100LGN00 \\	
			2023/02/21 & \makecell{LC09\_L1TP\_191035\_20230221\_20230309\_02\_T1} & \makecell{LC09\_L2SP\_191035\_20230221\_20230309\_02\_T1} & LC91910352023052LGN01 \\	
			2017/07/06 & \makecell{LC08\_L1TP\_191035\_20170706\_20200903\_02\_T1} & \makecell{LC08\_L2SP\_191035\_20170706\_20200903\_02\_T1} & LC81910352017187LGN00 \\	
			2017/04/01 & \makecell{LC08\_L1TP\_191035\_20170401\_20200904\_02\_T1} & \makecell{LC08\_L2SP\_191035\_20170401\_20200904\_02\_T1} & LC81910352017091LGN00 \\	
			2017/02/28 & \makecell{LC08\_L1TP\_191035\_20170228\_20200905\_02\_T1} & \makecell{LC08\_L2SP\_191035\_20170228\_20200905\_02\_T1} & LC81910352017059LGN00 \\	
			\midrule
			\multicolumn{4}{c}{\textbf{Gulf of Gabès region}}\\
			\midrule
			2023/07/15 & \makecell{LC09\_L1TP\_191036\_20230715\_20230715\_02\_T1} & \makecell{LC09\_L2SP\_191036\_20230715\_20230717\_02\_T1} & LC91910362023196LGN00 \\	
			2023/04/10 & \makecell{LC09\_L1TP\_191036\_20230410\_20230410\_02\_T1} & \makecell{LC09\_L2SP\_191036\_20230410\_20230412\_02\_T1} & LC91910362023100LGN00 \\	
			2023/02/21 & \makecell{LC09\_L1TP\_191036\_20230221\_20230309\_02\_T1} & \makecell{LC09\_L2SP\_191036\_20230221\_20230309\_02\_T1} & LC91910362023052LGN01 \\	
			2017/07/22 & \makecell{LC08\_L1TP\_191036\_20170722\_20200903\_02\_T1} & \makecell{LC08\_L2SP\_191036\_20170722\_20200903\_02\_T1} & LC81910362017203LGN00 \\
			2017/04/01 & \makecell{LC08\_L1TP\_191036\_20170401\_20200904\_02\_T1} & \makecell{LC08\_L2SP\_191036\_20170401\_20200904\_02\_T1} & LC81910362017091LGN00 \\
			2017/02/28 & \makecell{LC08\_L1TP\_191036\_20170228\_20200905\_02\_T1} & \makecell{LC08\_L2SP\_191036\_20170228\_20200905\_02\_T1} & LC81910362017059LGN00 \\
			\bottomrule
		\end{tabularx}
\end{adjustwidth}
\end{table} 

All images are selected on the three key control periods: winter with the highest water level (February), spring with transitional period and rains (April) and summer with the peak heat (July). The metadata in the images are listed in the Tables in Appendix \ref{AppA}. It is widely accepted that cloudiness and haze has a key impact on the image properties. Therefore, in this study we used the 10\% of cloud coverage for all images. Tested images included spectral bands in visible, Short-wave infrared (SWIR)-1, SWIR2 and Near-infrared (NIR) channels obtained from the OLI sensor. The topographic maps showing two study areas within the country and their enlarged fragments are performed with the Generic Mapping Tools (GMT) software, version  6.1.1 \cite{Wessel} using the existing scripting techniques reported earlier \cite{Lemenkova202212,Lemenkova202206}. 

%----------- subsection ----------->
\subsection{Methods}

The general scheme of the data and methods applied in this study is presented schematically in a flowchart in Figure \ref{fig02}. The Landsat 8-9 OLI/TIRS satellite images were processed with the Geographic Resources Analysis Support System (GRASS) Geographic Information System (GIS) software. Compared to the existing GIS methods, the GRASS GIS presents an advanced method of geographic visualization. Furthermore, the increase of the open spatial data arises a question of the effective tools for their processing. Since GIS have restricted functionality, modern approaches such as GRASS GIS present a powerful alternative through scripting, which presents a principally new paradigm for cartographic data processing and satellite image analysis. 

The image processing was performed on the MacBookAir Apple computer with chip Apple M1 which allowed to technically process large files with high amount of the required computer memory. First, the images were processed by scripts using clustering using 'i.cluster' module and continued with classification by the 'i.maxlik' module. The scenes were analysed for 10 recent years from 2013 to 2023 using the maximum-likelihood discriminant analysis classifier. Afterwards, the spatio-temporal changes of the water level and salinisation in the target salt lakes of Tunisia and its climate-environmental driving forces were analysed. The detection of the saline lakes is possible using spectral properties of the images that enable to distinct highly reflective cyan-coloured areas of the lakes discriminated against the beige-coloured sands and bare land. Thus, spectral properties of the RS data were evaluated by various GRASS GIS modules and scripting techniques with regards to the water/land/salt discrimination, and demonstrated robust approach to image analysis. 

Technically, using scripts optimises the use of the computer memory and computation resources. The snippets of GRASS GIS code are concatenated into longer scripts which perform the complete task of mapping. Such a partition of the workflow enables the machine to perform selected tasks of image processing in a distributed way. This decreases a possibility of errors and misclassification of the pixels and supports processing large datasets due to a significant level of automation. Motivated by powerful functionality of GRASS GIS and its modules for image processing, this study presents a case of processing remote sensing data by scripts. 

\begin{figure}[H]
\hspace{50pt}\makebox[0.90\linewidth][r]{%
	\centering
		\includegraphics[width=15.0 cm]{Fig_02.jpg}
}
\caption{Flowchart summarising the workflow. Software: R. Diagram source: author.}\label{fig02}
\end{figure}

\subsubsection{Data preprocessing}
First, the images were preprocessed by the 'i.landsat.toar' module of GRASS GIS using properties of image and information on sun elevation level to evaluate the reflectance in pixels and calculation of top-of-atmosphere radiance and temperature for all the Landsat OLI scenes. Hence, the images were corrected and added into GRASS environment using 'r.import' module. After the import procedure, the images with all the multispectral bands were grouped in the working directory using 'i.group' module, which is used to select necessary bands, and their data were checked using 'g.list' module. The 'i.group' module of the GRASS GIS was used for processing the structures of the spatial data and grouping them into classes using the k-means algorithm. 

The metadata of the images were checked using relevant modules of the GRASS GIS. In particular, the techniques of image processing in GRASS GIS include the auxiliary functions of 'g.list rast' and 'r.info' acting on imported raster bands and showing the identified parameters of images as attributes and a given set of identifiers of technical specifications of the multispectral images: resolution, minimal and maximal values, coordinate system. After image preprocessing, the 'i.group' module was used for grouping the bands before the classification of the Landsat images. This approach partitions the image using the embedded algorithm of discriminating the pixels on the images that represent various land cover types in target area using their spectral reflectances. Afterwards, they were clustered using the k-means algorithm embedded in the GRASS GIS which was implemented by the 'i.cluster' module. 

%----------- subsection ----------->
\subsubsection{Data clustering}

During clustering, a 'signature file' was created which was used in the next step to sort pixels into the land cover classes using 'i.cluster' module of the GRASS GIS which creates cluster and covariance matrices. All classes contained pixels with approximately the same level of spectral reflectance, which was automatically evaluated by the machine based analysis of the GRASS GIS but different number of pixels within each class and thus different extent of landscape categories, which enabled us to evaluate the patches and their distribution within the scenes under the effects of various seasons on development of salinisation in sabhkas. Besides, the images located in the Gulf of Gabès contained more sandy desert areas compared to the northern segment of the study areas in the Gulf of Hammamet which allowed to compare the effects from geographic distribution of the saline lakes with higher coverage of mixed vegetation in the more temperate climate of north compared to the southern region which is more impacted by the arid climate of Sahara.

\begin{figure}[H]
\hspace{50pt}\makebox[0.90\linewidth][r]{%
	\begin{subfigure}[b]{.40\textwidth}
		\centering
			\includegraphics[width=0.87\textwidth]{Fig_03a.jpg}
		\caption{28 February 2017}
	\end{subfigure}%
	\begin{subfigure}[b]{.40\textwidth}
		\centering
		\includegraphics[width=0.87\textwidth]{Fig_03b.jpg}
		\caption{1 April 2017}
	\end{subfigure}%
	\begin{subfigure}[b]{.40\textwidth}
		\centering
		\includegraphics[width=0.87\textwidth]{Fig_03c.jpg}
		\caption{6 July 2017}
	\end{subfigure}%
}\\
\vfill \vspace{1mm}
\hspace{50pt}\makebox[0.90\linewidth][r]{%
	\begin{subfigure}[b]{.40\textwidth}
		\centering
			\includegraphics[width=0.87\textwidth]{Fig_03d.jpg}
		\caption{21 February 2023}
	\end{subfigure}%
	\begin{subfigure}[b]{.40\textwidth}
		\centering
		\includegraphics[width=0.87\textwidth]{Fig_03e.jpg}
		\caption{10 April 2023}
	\end{subfigure}%
	\begin{subfigure}[b]{.40\textwidth}
		\centering
		\includegraphics[width=0.87\textwidth]{Fig_03f.jpg}
		\caption{15 July 2023}
	\end{subfigure}%
}
\caption{Landsat images on the Gulf of Hammamet region (northern segment) showing the increase of salt (bright cyan colour) from winter to summer months in sabkhas Sebkhet de Sidi el Hani (Sousse Governorate), Sebkha de Moknine (Mahdia Governorate) and Sebkhet El Rharra (Sfax Governorate): \textbf{(a)}: 28 February 2017; \textbf{(b)}: 1 April 2017; \textbf{(c)}: 6 July 2017; \textbf{(d)}: 21 February 2023; \textbf{(e)}: 10 April 2023; \textbf{(f)}: 15 July 2023.}\label{fig03}
\end{figure}

%----------- subsection ----------->
\subsubsection{Image classification}

Theoretically, the essential  principle of the land cover type classification is based on the fundamental properties of RS which consists in measurement of spectral reflectance as a reflection of light from the Earth's surface. Spectral reflectance signatures of various land cover types (land, water, salt soils, vegetation, urban areas, etc) differs significantly. This is because of the different emitted radiation from these these objects which varies significantly. The Landsat 8-9 OLI/TIRS sensor continuously measures spectral reflectance of the objects on the Earth which are identified as various land cover types. Hence, the performed measurements indicate various different bodies and objects on the Earth using their spectral reflectance visible on the spaceborne satellite images in selected band combinations. Consequently, since measured land cover types varies, it enables to perform classification of the land cover types using RS data processing, following by cartographic visualization and interpretation. 

Practically, the information on spectral reflectance is used for classification procedure. Thus, after clustering, the images were processed using special module of the GRASS GIS which performs unsupervised automatic classification of the satellite images. The advanced methods of GRASS GIS present effective instruments for classification -- the 'i.maxlik', developed by M. Shapiro and T. Wen, and supported with additional module by M. Nartiss. The use of the 'i.maxlik' only required to specify the parameters of signature file and related information on groups and subgroups of spectral bands during data partition, and to assign pixels into each land cover class according to the records of the spectral reflectance in binary format. The spectral reflectance of wavelengths which develops proportionally with internal gain in strength in brightness of land cover types identified from space by the sensor. The technical characteristics of this module include two steps: the preceding clustering and following classification. The module is designed to be used for the multispectral images, such as Landsat, using spectral signatures of the pixels. 

\begin{figure}[H]
\hspace{50pt}\makebox[0.90\linewidth][r]{%
	\begin{subfigure}[b]{.40\textwidth}
		\centering
			\includegraphics[width=0.87\textwidth]{Fig_04a.jpg}
		\caption{28 February 2017}
	\end{subfigure}%
	\begin{subfigure}[b]{.40\textwidth}
		\centering
		\includegraphics[width=0.87\textwidth]{Fig_04b.jpg}
		\caption{1 April 2017}
	\end{subfigure}%
	\begin{subfigure}[b]{.40\textwidth}
		\centering
		\includegraphics[width=0.87\textwidth]{Fig_04c.jpg}
		\caption{22 July 2017}
	\end{subfigure}%
}\\
\vfill \vspace{1mm}
\hspace{50pt}\makebox[0.90\linewidth][r]{%
	\begin{subfigure}[b]{.40\textwidth}
		\centering
			\includegraphics[width=0.87\textwidth]{Fig_04d.jpg}
		\caption{21 February 2023}
	\end{subfigure}%
	\begin{subfigure}[b]{.40\textwidth}
		\centering
		\includegraphics[width=0.87\textwidth]{Fig_04e.jpg}
		\caption{10 April 2023}
	\end{subfigure}%
	\begin{subfigure}[b]{.40\textwidth}
		\centering
		\includegraphics[width=0.87\textwidth]{Fig_04f.jpg}
		\caption{15 July 2023}
	\end{subfigure}%
}
\caption{Landsat images on the Gulf of Gabès region (southern segment) showing the increase of salt (bright cyan colour) from winter to summer months in sabkhas Sebkhet en Noual (Sfax Governorate): \textbf{(a)}: 28 February 2017; \textbf{(b)}: 1 April 2017; \textbf{(c)}: 6 July 2017; \textbf{(d)}: 21 February 2023; \textbf{(e)}: 10 April 2023; \textbf{(f)}: 15 July 2023.}\label{fig04}
\end{figure}

The classification procedure was performed with different images (6 scenes for the study area-1 in the Gulf of Gabès, and 6 images for the study area-2 for the Gulf of Hammamet). The aim of the classification to identify cells in the classified image and to group then into classes according to the similarity of spectral reflectance in each given group (class) through estimated spectral reflectance value during automatic image analysis performed by the GRASS GIS algorithm. As for comparison of several images taken on various seasons and various years with 6 year interval, the main idea was to compare the effects of climate, the impact from the Sahara and seasonal meteorological fluctuations on the development of salinisation in sabkhas of Tunisia using comparative analysis. Therefore, the experiments with images were implemented with different calendar seasons and in different regions (central/southern) of Tunisia. 

%----------- subsection ----------->

 \subsubsection{Accuracy analysis}

The images were evaluated for accuracy and pixels with high level of rejection were filtered out using 'r.mapcalc' module of GRASS GIS. The rejection probability values, i.e., the pixel classification confidence levels were computed and visualised and a series of maps. The error matrix and  convergence level of classified pixels were measured using 'i.cluster' module that includes accuracy assessment of classification results in tabular format. After that, the classification was performed by the and the results of the clustering were exported as a series of matrices in CSV format for each images, and the results are reported in the tables in the Appendix \ref{AppB} of this study. The classified satellite scenes were verified and class separability matrix was computed for each of the 12 images. The accuracy of the classification was measured and summarised in the tables with a convergence level of classified pixels above 98\% of points stable. Such approach to analyse the accuracy of image processing provided data to evaluate the obtained values of land cover classes at each image of the processed time series.
 
 %----------- subsection ----------->
 \subsubsection{Mapping}

The environmental properties of the landscapes were evaluated with cartographic methods after image processing using visualised maps. The aim is to analyse the seasonal effects of climate (increase of temperature and evaporation) on the behaviour of salt lakes. The complete cartographic workflow used in this study included a sequence of the following modules of GRASS GIS:

\begin{itemize}
	\item 'r.import' using for data import via the embedded GDAL library; the extent and resolution were adjusted to the region of the Landsat scenes;
	\item 'g.list' for listing imported raster data type by the search pattern of GRASS GIS; 
	\item 'g.region' for defining the computational extent of the region to match the scenes; 
	\item 'i.cluster' for clustering of the raster Landsat images using k-means algorithm; 
	\item 'i.maxlik' for unsupervised image classification using maximum likelihood method; 
	\item 'd.mon' for display visualisation and cartographic plotting; 
	\item 'r.colors' for adjusting the colour palette according to the land cover types; 
	\item 'd.rast' for mapping the image in the active graphics frame; 
	\item 'd.legend' for adding the colour legend with textual notations; 
	\item 'd.out.file' for data export to the conventional formats (JPG). 
\end{itemize}

These GRASS GIS modules have been incorporated to perform various steps of cartographic tasks and image processing workflow by using the GRASS GIS console. The land cover types were identified continuously across the images using a sequence of k-means clustering and maximum likelihood classification approaches, and converted to the maps which indicate variations in land cover types over Tunisia during a period of 2017 to 2023 for two image datasets covering the regions on the Gulf of Hammamet and the Gulf of Gabès. The data obtained from the GRASS GIS based image processing workflow were evaluated for environmental analysis.

%----------- section ----------->
\section{Results and Discussion}

The results of the image processing approach of the GRASS GIS demonstrated in this paper show that the summer period speeds up soil salinisation and mineralisation in sabkhas of Tunisia, as demonstrated through mapped series of images on several seasons with increased summer heat, and vice versa. Then, using the information received from spectral reflectance of the pixels on the images and heat development in Sahara during summer months we were able to find the highest degree of salinisation for each lake in various years. Figure \ref{fig05} demonstrates the results of the image classification showing the behaviour of saline lakes over seasonal between 2017 and 2023 in the region of Gulf of Hammamet (Sebkhet de Sidi el Hani, Sebkha de Moknine and Sebkhet El Rharra). 

\begin{figure}[H]
\hspace{50pt}\makebox[0.90\linewidth][r]{%
	\begin{subfigure}[b]{.40\textwidth}
		\centering
			\includegraphics[width=0.95\textwidth]{Fig_05a.jpg}
		\caption{28 February 2017}
	\end{subfigure}%
	\begin{subfigure}[b]{.40\textwidth}
		\centering
		\includegraphics[width=0.95\textwidth]{Fig_05b.jpg}
		\caption{1 April 2017}
	\end{subfigure}%
	\begin{subfigure}[b]{.40\textwidth}
		\centering
		\includegraphics[width=0.95\textwidth]{Fig_05c.jpg}
		\caption{6 July 2017}
	\end{subfigure}%
}\\
\vfill \vspace{1mm}
\hspace{50pt}\makebox[0.90\linewidth][r]{%
	\begin{subfigure}[b]{.40\textwidth}
		\centering
			\includegraphics[width=0.95\textwidth]{Fig_05d.jpg}
		\caption{21 February 2023}
	\end{subfigure}%
	\begin{subfigure}[b]{.40\textwidth}
		\centering
		\includegraphics[width=0.95\textwidth]{Fig_05e.jpg}
		\caption{10 April 2023}
	\end{subfigure}%
	\begin{subfigure}[b]{.40\textwidth}
		\centering
		\includegraphics[width=0.95\textwidth]{Fig_05f.jpg}
		\caption{15 July 2023}
	\end{subfigure}%
}
\caption{Classification results showing land cover classes based on the processed Landsat images showing northern study area (Gulf of Hammamet): Sebkhet de Sidi el Hani (Sousse Governorate), Sebkha de Moknine (Mahdia Governorate) and Sebkhet El Rharra (Sfax Governorate): \textbf{(a)}: 28 February 2017; \textbf{(b)}: 1 April 2017; \textbf{(c)}: 6 July 2017; \textbf{(d)}: 21 February 2023; \textbf{(e)}: 10 April 2023; \textbf{(f)}: 15 July 2023.}\label{fig05}
\end{figure}

Figure \ref{fig06} demonstrates the outputs of the image classification showing variations in land cover types and saline lakes over seasonal between 2017 and 2023 in the region of Gulf of Gabès (Sebkhet en Noual). The land cover types were identified on the images using information adopted from and modified using the existing classification scheme on land cover types in Tunisia \cite{plants10102031} and data from Food and Agriculture Organization (FAO) \cite{FAO}. The existing land cover types for the country-level scale were modified according to the extent of the images covering study area and included the following types: 1) forest; 2) grasslands and rangelands; 3) sand and desert; 4) cultivated agriculture plantations and arable lands (olive tree, orchards, vine, palm groves); 5) urban areas and artificial land; 6) wasteland; 7) wetlands and water; 8) bare land and soil; 9) mixed vegetation; 10) saline soil. 

Changed seasons enabled to monitor and interpret how the salinisation process activate the accumulation of minerals in the salt lakes with changed colour of sabkhas. Thus, during the summer period, the structure of the crystals in the soil sediments in sabkhas was solidified by the increased evaporation as a result of the increased temperatures. In turn, the seasonal effects of increased temperature and lack of precipitation in the southern region (Figure \ref{fig06}) are reflected in the increased areas of sand and declined vegetation, as shown in the maps. Changes in salt lakes were determined at different years during recent decade by integrating Landsat images taken in winter/spring/summer months for comparison of periods when water inflow into Saharan lakes is high and low, respectively. Note that the decline in vegetation areas is clear on the images of 2017 from April to July showing sabkhas Sebkhet de Sidi el Hani, Sebkha de Moknine  and Sebkhet El Rharra, Figure \ref{fig06}. 

Compared to the difference between the February and April months, the decline in vegetation and increase in salinisation is more prominent, because it indicates the increase in summer heat temperatures which was contrasting compared to the spring months. In Figure \ref{fig06}, it is seen that the results are quite different for visualised series of maps for various seasons, because salt lakes contain different amounts of water with regard to salt, i.e., the classified images revealed various water/salt ratio in the sabkhas along with the development of heat during summer months and effects from Sahara desert in the southern study region. The classification of pixel confidence levels with rejection probability values indicating the accuracy of classification are presented in Figure \ref{fig07} for the northern region of the Gulf of Hammamet, and in Figure \ref{fig08} for the southern region of the Gulf of Gabès, respectively.

\begin{figure}[H]
\hspace{50pt}\makebox[0.90\linewidth][r]{%
	\begin{subfigure}[b]{.40\textwidth}
		\centering
			\includegraphics[width=0.95\textwidth]{Fig_06a.jpg}
		\caption{28 February 2017}
	\end{subfigure}%
	\begin{subfigure}[b]{.40\textwidth}
		\centering
		\includegraphics[width=0.95\textwidth]{Fig_06b.jpg}
		\caption{1 April 2017}
	\end{subfigure}%
	\begin{subfigure}[b]{.40\textwidth}
		\centering
		\includegraphics[width=0.95\textwidth]{Fig_06c.jpg}
		\caption{6 July 2017}
	\end{subfigure}%
}\\
\vfill \vspace{1mm}
\hspace{50pt}\makebox[0.90\linewidth][r]{%
	\begin{subfigure}[b]{.40\textwidth}
		\centering
			\includegraphics[width=0.95\textwidth]{Fig_06d.jpg}
		\caption{21 February 2023}
	\end{subfigure}%
	\begin{subfigure}[b]{.40\textwidth}
		\centering
		\includegraphics[width=0.95\textwidth]{Fig_06e.jpg}
		\caption{10 April 2023}
	\end{subfigure}%
	\begin{subfigure}[b]{.40\textwidth}
		\centering
		\includegraphics[width=0.95\textwidth]{Fig_06f.jpg}
		\caption{15 July 2023}
	\end{subfigure}%
}
\caption{Classification results showing land cover classes based on the processed Landsat images showing southern study area (Gulf of Gabès), Sebkhet en Noual (Sfax Governorate): \textbf{(a)}: 28 February 2017; \textbf{(b)}: 1 April 2017; \textbf{(c)}: 6 July 2017; \textbf{(d)}: 21 February 2023; \textbf{(e)}: 10 April 2023; \textbf{(f)}: 15 July 2023.}\label{fig06}
\end{figure} 

Varied temperature in seasonal climate periods has effects on the behaviour of lakes and salinisation of sabkhas which varies for different lakes. This also serves as indirect indicator of groundwater inflow and impacts of coastal climate on the eastern sabkhas, for instance for the case of Sebkha de Moknine which is located near the Bekalta coastal town. This lake experiences the least difference in the salinisation and was filled with water during the 2017, while in 2023 it experienced salinisation only during summer. In contrast, Sebkhet El Rharra experiences clear fluctuations in the salinisation during each year being regularly filled with water in winter and covered by salt crust during July. Such changes were detected both for year 2017 and for year 2023.

The extends of agricultural areas are much more similar in Figure \ref{fig06} for the Gulf of Gabès area, indicating that the plantation areas in southern regions are rather similar for different landscape spots compared to the northern segment (Figure \ref{fig05}). This is also seen in Figure \ref{fig06} for 2017 and 2023 years on the month of April. The use of comparative analysis of several images showing dynamics of land cover types as descriptors to reveal the development of aridification and desertification in southern regions closer to Sahara in study areas demonstrated RS approach to be a straightforward yet efficacious strategy to extract contextual information regarding land cover properties of landscapes. Thus, although RS data cannot directly detect landscape dynamics since this is dependent on the comparison of land cover properties for several scenes taken in different periods, nevertheless, processing of RS data enables to approximate the extent of landscape patches from classified images and evaluate water content in salt lakes as indirect indicators of salinisation during summer heat using image processing with the GIS methods. 

\begin{figure}[H]
\hspace{50pt}\makebox[0.90\linewidth][r]{%
	\begin{subfigure}[b]{.40\textwidth}
		\centering
			\includegraphics[width=0.95\textwidth]{Fig_07a.jpg}
		\caption{28 February 2017}
	\end{subfigure}%
	\begin{subfigure}[b]{.40\textwidth}
		\centering
		\includegraphics[width=0.95\textwidth]{Fig_07b.jpg}
		\caption{1 April 2017}
	\end{subfigure}%
	\begin{subfigure}[b]{.40\textwidth}
		\centering
		\includegraphics[width=0.95\textwidth]{Fig_07c.jpg}
		\caption{6 July 2017}
	\end{subfigure}%
}\\
\vfill \vspace{1mm}
\hspace{50pt}\makebox[0.90\linewidth][r]{%
	\begin{subfigure}[b]{.40\textwidth}
		\centering
			\includegraphics[width=0.95\textwidth]{Fig_07d.jpg}
		\caption{21 February 2023}
	\end{subfigure}%
	\begin{subfigure}[b]{.40\textwidth}
		\centering
		\includegraphics[width=0.95\textwidth]{Fig_07e.jpg}
		\caption{10 April 2023}
	\end{subfigure}%
	\begin{subfigure}[b]{.40\textwidth}
		\centering
		\includegraphics[width=0.95\textwidth]{Fig_07f.jpg}
		\caption{15 July 2023}
	\end{subfigure}%
}
\caption{Classification of pixel confidence levels with rejection probability values for the northern study area (Gulf of Hammamet): Sebkhet de Sidi el Hani (Sousse Governorate), Sebkha de Moknine (Mahdia Governorate) and Sebkhet El Rharra (Sfax Governorate): \textbf{(a)}: 28 February 2017; \textbf{(b)}: 1 April 2017; \textbf{(c)}: 6 July 2017; \textbf{(d)}: 21 February 2023; \textbf{(e)}: 10 April 2023; \textbf{(f)}: 15 July 2023.}\label{fig07}
\end{figure}

A time series analysis of the satellite images demonstrated in this paper characterises gradual changes in lakes over calendar periods and provides a robust view in environmental changes. Here, the obtained results presented a time series of processed satellite images that showed that the increase in salt lakes and desertification of surrounding areas is more distinct in the southern segment which showed that the Sebkhet en Noual is prone to salinisation even during the winter months, while the northern sabkhas Sebkhet de Sidi el Hani , Sebkha de Moknine and Sebkhet El Rharra show water-filled basins during winter. Furthermore, it demonstrated that the regions of Sfax are more prone to aridification compared to the areas of Mahdia and Sousse. 

Changes in water level in salt lakes were detected using RS data. The presented results can serve for environmental monitoring of Tunisian landscapes and analysis of environmental consequences of climate effects in central and north Tunisia. Although the results of image processing that analyse the environmental landscape properties and changes in the land cover types have been described earlier in various works on Tunisia \cite{Touhami,Khelifa,Hammami,Bounouh} providing a general pattern of land cover change by various GIS, the results proposed here are obtained using recent Landsat OLI/TIRS images processed by the GRASS GIS scripts. For instance, earlier works use the Landsat MSS/4 and SPOT HRV/3 images \cite{Park}, Landsat and SRTM processed by ArcGIS and ENVI commercial software \cite{Souissi,Afrasinei} or a combination of ETM+ and OLI 8 sensors of the Landsat products \cite{Mezned}. 

\begin{figure}[H]
\hspace{50pt}\makebox[0.90\linewidth][r]{%
	\begin{subfigure}[b]{.40\textwidth}
		\centering
			\includegraphics[width=0.95\textwidth]{Fig_08a.jpg}
		\caption{28 February 2017}
	\end{subfigure}%
	\begin{subfigure}[b]{.40\textwidth}
		\centering
		\includegraphics[width=0.95\textwidth]{Fig_08b.jpg}
		\caption{1 April 2017}
	\end{subfigure}%
	\begin{subfigure}[b]{.40\textwidth}
		\centering
		\includegraphics[width=0.95\textwidth]{Fig_08c.jpg}
		\caption{6 July 2017}
	\end{subfigure}%
}\\
\vfill \vspace{1mm}
\hspace{50pt}\makebox[0.90\linewidth][r]{%
	\begin{subfigure}[b]{.40\textwidth}
		\centering
			\includegraphics[width=0.95\textwidth]{Fig_08d.jpg}
		\caption{21 February 2023}
	\end{subfigure}%
	\begin{subfigure}[b]{.40\textwidth}
		\centering
		\includegraphics[width=0.95\textwidth]{Fig_08e.jpg}
		\caption{10 April 2023}
	\end{subfigure}%
	\begin{subfigure}[b]{.40\textwidth}
		\centering
		\includegraphics[width=0.95\textwidth]{Fig_08f.jpg}
		\caption{15 July 2023}
	\end{subfigure}%
}
\caption{Classification of pixel confidence levels with rejection probability values for the southern study area (Gulf of Gabès): Sebkhet en Noual (Sfax Governorate: \textbf{(a)}: 28 February 2017; \textbf{(b)}: 1 April 2017; \textbf{(c)}: 6 July 2017; \textbf{(d)}: 21 February 2023; \textbf{(e)}: 10 April 2023; \textbf{(f)}: 15 July 2023.}\label{fig08}
\end{figure}

The reaction of materials sand salinisation processes during high evaporation soil results in changed structure of pores of soil with significantly reduced water content in soil material during periodic occurrences of a deep brine layer on the bottom of the lakes. Such behaviour of soil as lacustrine sediment layer demonstrates the response to climate and meteorological processes so that the considerable fluctuations in water level altered the soil structure in lacustrine sediments which is reflected as increased salinity and decreased water level, as evaluated during winter/spring/summer periods on satellite images. Moreover, the degree in salinisation of sabkhas well correlates with the location of each lake with southern lakes almost completely covered by salt even during summer, while northern lakes demonstrate fluctuation. 

Similar to existing studies \cite{BAYATI2021126032,SHENG2016129,SCHRODER20221710}, the salinisation of lacustrine soil evaluated using RS data, this study showed the highest level of changes in the July month, and the most contrasting behaviour of brine accumulation in the lake of Sebkhet de Sidi el Hani (Sousse region) which has two connected sub-basins with different morphology hat might have effects on salinisation of the lake. The advanced scripting cartographic approach of GRASS GIS was introduced to evaluate the lacustrine properties of soil around the sabkhas in north Tunisia using RS data processing. It is well known that proprietary GIS and RS software are expensive since their development is costly (e.g., ArcGIS or Erdas Imagine). Moreover, very high resolution satellite images are not always easily accessible since they are provided on the commercial basis (e.g., SPOT, IKONOS). 

Therefore, a special advantage of this work is the use of the open source Landsat OLI/TIRS data and a free GRASS GIS software. This makes the present work repeatable, which can be continued in similar relevant studies in the future projects. The presented open source solutions on the optimal selection of data and tools can be supportive for environmental monitoring  of Tunisia. Hence, a series of the satellite images were processed with various combination of GRASS GIS modules and compared with the images covering identical areas on earlier periods, which was experimentally used for environmental monitoring of sabkhas in Tunisia. The performed experiments on RS data analysis supported monitoring of fluctuations in water level of salt lakes of Tunisia through the analysed time series analysis of satellite images. 

%----------- section ----------->
\section{Conclusions}

\subsection{Summary}

GRASS GIS software has been successfully applied to identify different stages in the landscape change in sabkhas of north Tunisia (with a special focus on several lakes -- Sebkhet de Sidi el Hani, Sebkha de Moknine, Sebkhet El Rharra and Sebkhet en Noual). The study also evaluated the effects of climate on seasonal changes in the lacustrine landscapes, as indicated by development of land cover types over the period of 2013 to 2023. Besides, the experiments served as a means of evaluating the climate effects as accelerators of salinisation development in sediments of lakes. The presented results confirm that GRASS GIS can be effectively used to processing and classification of RS imagery using programming approach that ensures automation and precision and accuracy of data processing. 

\subsection{Recommendations}

For future similar works on RS data analysis of salt lakes, we provide below the following recommendations and technical suggestions for mapping and image analysis:

\begin{enumerate}
	\item Make the initial study similar to the image processing workflow in this report, but also use other GRASS GIS modules applied for processing the target dates. For example, a useful approach is segmentation by 'i.segment'.
	\item Due to the  dynamics of plant phenology in different months, the period of image capture should be identical to evaluate changes for each image. For example, a correct comparison can be performed using images taken for a specific month (e.g., January) in a sequence of several years, or several identical months for two years, as presented in this study.
	\item Cloudiness vary a lot in different images. It is necessary to use a cloud-free series of images as an input dataset. While in this study, 10\% of cloud coverage was applied, the modified parameters optimised with different cloud content can be tested and applied with decreased cloud coverage. In this way, the use of 'i.landsat.acca' module of the GRSAS GIS enables Automatic Cloud Cover Assessment (ACCA) for a quick evaluation of image quality and suitability.
	\item Process other dataset in addition to the Landsat, e.g., Advanced Spaceborne Thermal Emission and Reflection Radiometer (ASTER) images. ASTER data can be proprocced using special modules of 'i.aster.toar', similar to the 'i.landsat.toar'. The main advantage of ASTER is a higher resolution--a pixel size of 15 m compared to the 30-m Landsat. 
	\item For mapping, take satellite images from the USGS on diverse dates to evaluate spatio-temporal dynamics and perform mapping on images captured on different seasonal periods to compare the difference in salinisation between the lakes.
	\item The  commercial GIS currently available on the market typically have license. Therefore, one can take open source GIS software such as QGIS or GRASS GIS, which have effective tools and extended functionality which are used for for RS and cartographic workload.
	\item Evaluate the integral effects from various factors that affect the landcover changes, i.e., the heat from the increased temperature, decreased precipitation and high evaporation lead to desertification of landscapes. 
	\item The access to the groundwaters and geographic location of the landscapes affect the behaviour of the vegetation. For instance, closeness to the desert or coastal areas affect differently vegetation patterns of the landscapes. 
	\item For analysis of vegetation health and vigour, the computation of the NDVI or other vegetation indices is a useful means of environmental assessment. In GRASS GIS, vegetation indices can be automatically calculated  by the software using the embedded formulae adjusted to the bands of the relevant satellite images. For example, NIR and Red channels correspond to the bands 5 and 4 for the Landsat 8-9 OLI/TIRS, while bands 4 and 3 for older Landsat 7 ETM+ images.
	\item Changes in land cover types within certain limits indicate that there is dynamics in landscapes as a response to the environmental and climate effects and processes. Moreover, it may indicate the cumulative effects of several factors on soil-vegetation conditions, including both climate-environmental processes and anthropogenic activities.
\end{enumerate}

\subsection{Concluding remarks}

This paper reports the results obtained from experimental investigations on satellite images Landsat 8-9 OLI/TIRS collected from the USGS. The tests on image processing were designed to evaluate the changes in several saline lakes of north Tunisia in various seasons in 2013 to 2023 and variations in water extent as indicators of salinisation. Specifically, we evaluated the proposed framework on RS data analysis using GRASS GIS software and a series of the satellite images from the USGS repository. We also proposed future directions in similar works to develop robust scheme on RS data processing tests in similar studies where cartographic approach to image analysis is applied in environmental monitoring projects aimed at evaluation of climate effects on landscape changes. In this work we also discussed the ways in which the scripting approach of cartographic data processing contrasts with traditional GIS methods of mapping that use Graphical User Interface (GUI), and complement to RS data processing for evaluating the connections between the climate effects and salinisation of sabkhas in Tunisia. 

We demonstrated the advantages of the GRASS GIS as a quantitative, robust and advanced cartographic tool for evaluating environmental properties based on image analysis. A workflow has been developed that uses a series of high-resolution Landsat 8-9 OLI/TIRS images with various combinations of GRASS GIS modules, evaluated to analyse the properties of lacustrine landscapes  of north Tunisia as a function of summer heat in Sahara and increase in evaporation during the period of 10 years (2013-2023). The approach of GRASS GIS demonstrated to be a more advanced way of image processing compared to traditional GIS software due to scripts that enables automation of workflow repeatability of the process. Thus, as showed in this study, the use of several module of GRASS GIS ensures multifold manipulations with images, such as in 'i.segment', 'i.maxlik' and other sub-tests that where applied for various steps of multispectral image analysis. 

A scripting approach of the GRASS GIS enabled us to map and analyse the rate of salinisation in sabkhas of Tunisia in real time in various seasons from 2013 to 2023 and to obtain the information from the processed images representing the land cover types in the surroundings. In turn, RS datasets enabled to model the behavior of landscapes in spring, summer and winter periods in north Tunisia which was used to compare the climate effects from rise in temperature on the evaporation of water in salt lakes, as directly related to heat increase during summer period in Sahara. The results demonstrated variations in water level in several sabkha and show technical effectiveness of the GRASS GIS in cartographic processing of the RS data for environmental tasks. Correlations between the increase in heat in various years and high salinisation of sabkhas were detected. The experiments on image processing reported in this article have investigated the salinisation behaviour in saline lakes with thermal variations using analysis of Landsat bands in multispectral channels. 

Due to the effectiveness and applicability of the performed experiments on image analysis by GRASS GIS -- flexibility of scripting approach and a wide variety of modules tuned for diverse tasks of image processing using its modular approach -- we expect the proposed method to be able to deal with other satellite imagery that cover diverse regions of north Africa to study development of lacustrine landscapes with effects from climate and fluctuations in temperature on Saharan ecosystems. Various sabkhas of north Tunisia with diverse surrounding soil types and salinisation/mineralisation degree were evaluated to generalise the proposed workflow of salt lake monitoring to a higher dimensional workload in the larger lakes and different world ecosystems with similar salt lakes. 

As demonstrated in this study, the RS data from the USGS EarthExplorer repository can be used to analyse the variations in salt lakes. It enables to derive information using data on spectral reflectance and brightness of pixels which strongly correlates with the distribution of various land cover types. In such a way, remote sensing data processing enables to detect patches in complex diversified lacustrine landscapes of north Sahara characterised with heterogeneous environment in contrasting climate of desert. This paper contributes to the environmental monitoring of Sahara, analysis of salt lake variations and regional studies of Tunisia, north Africa.

%----------- section ----------->

\funding{The publication was funded by the Editorial Office of \emph{Environments}, Multidisciplinary Digital Publishing Institute (MDPI), by providing 100\% discount for the APC of this manuscript.}

\institutionalreview{Not applicable.}

\informedconsent{Not applicable.}

\dataavailability{Not applicable} 

\acknowledgments{The author thanks the reviewers for reading and review of this manuscript.}

\conflictsofinterest{The author declares no conflict of interest.} 

\abbreviations{Abbreviations}{
The following abbreviations are used in this manuscript:\\
\noindent 
\begin{tabular}{@{}ll}
ACCA & Automatic Cloud Cover Assessment \\
ASTER & Advanced Spaceborne Thermal Emission and Reflection Radiometer \\
FAO & Food and Agriculture Organization \\
GDAL & Geospatial Data Abstraction Library \\
GMT & Generic Mapping Tools \\
GIS & Geographic Information System \\
GUI & Graphical User Interface \\
GRASS & Geographic Resources Analysis Support System \\
Landsat OLI/TIRS & Landsat Operational Land Imager and Thermal Infrared Sensor \\
MODIS & Moderate Resolution Imaging Spectroradiometer \\ 
NIR & Near-infrared \\
RS & Remote Sensing \\
SPOT & Satellite Pour l’Observation de la Terre \\
SWIR & Short-wave infrared (SWIR)-1 \\
USGS & United States Geological Survey \\
\end{tabular}
}

\appendixtitles{yes} 
\appendixstart
\appendix
\section[\appendixname~\thesection]{Metadata}\label{AppA}
\subsection[\appendixname~\thesubsection]{Gulf of Hammamet area}

\begin{table}[H]
\scriptsize
\caption{Metadata of Landsat OLI/TIRS satellite images for the Gulf of Hammamet area (year 2023).\label{tabApp}}
\begin{adjustwidth}{-\extralength}{0.5cm}
\centering
\begin{tabularx}{\fulllength}{|l|l|l|l}
    \toprule
        Data Set Attribute & Attribute Value & Attribute Value & Attribute Value \\ \hline
        Landsat Scene Identifier & LC91910352023196LGN00 & LC91910352023100LGN00 & LC91910352023052LGN01 \\ \hline
        Date Acquired & 2023/07/15 & 2023/04/10 & 2023/02/21 \\ \hline
        Collection Category & T1 & T1 & T1 \\ \hline
        Collection Number & 2 & 2 & 2 \\ \hline
        WRS Path & 191 & 191 & 191 \\ \hline
        WRS Row & 35 & 35 & 35 \\ \hline
        Target WRS Path & 191 & 191 & 191 \\ \hline
        Target WRS Row & 35 & 35 & 35 \\ \hline
        Nadir/Off Nadir & NADIR & NADIR & NADIR \\ \hline
        Roll Angle & 0.000 & 0.000 & 0.000 \\ \hline
        Date Product Generated L2 & 2023/07/17 & 2023/04/12 & 2023/03/09 \\ \hline
        Date Product Generated L1 & 2023/07/15 & 2023/04/10 & 2023/03/09 \\ \hline
        Start Time & 2023-07-15 09:54:18 & 2023-04-10 09:54:36 & 2023-02-21 09:55:02 \\ \hline
        Stop Time & 2023-07-15 09:54:50 & 2023-04-10 09:55:08 & 2023-02-21 09:55:33 \\ \hline
        Station Identifier & LGN & LGN & LGN \\ \hline
        Day/Night Indicator & DAY & DAY & DAY \\ \hline
        Land Cloud Cover & 0.01 & 0.59 & 1.20 \\ \hline
        Scene Cloud Cover L1 & 4.54 & 0.44 & 1.11 \\ \hline
        Ground Control Points Model & 1151 & 1119 & 1100 \\ \hline
        Ground Control Points Version & 5 & 5 & 5 \\ \hline
        Geometric RMSE Model & 5.646 & 5.848 & 6.402 \\ \hline
        Geometric RMSE Model X & 3.860 & 4.125 & 4.479 \\ \hline
        Geometric RMSE Model Y & 4.120 & 4.145 & 4.575 \\ \hline
        Processing Software Version & LPGS\_16.3.0 & LPGS\_16.2.0 & LPGS\_16.2.0 \\ \hline
        Sun Elevation L0RA & 65.60184606 & 55.81133510 & 38.05813927 \\ \hline
        Sun Azimuth L0RA & 120.36949136 & 140.16952296 & 149.12922810 \\ \hline
        TIRS SSM Model & N/A & N/A & N/A \\ \hline
        Data Type L2 & OLI\_TIRS\_L2SP & OLI\_TIRS\_L2SP & OLI\_TIRS\_L2SP \\ \hline
        Sensor Identifier & OLI\_TIRS & OLI\_TIRS & OLI\_TIRS \\ \hline
        Satellite & 9 & 9 & 9 \\ \hline
        Product Map Projection L1 & UTM & UTM & UTM \\ \hline
        UTM Zone & 32 & 32 & 32 \\ \hline
        Datum & WGS84 & WGS84 & WGS84 \\ \hline
        Ellipsoid & WGS84 & WGS84 & WGS84 \\ \hline
        Scene Center Lat DMS & "36°02'35.30""N" & "36°02'35.20""N" & "36°02'35.77""N" \\ \hline
        Scene Center Long DMS & "10°19'26.26""E" & "10°19'06.35""E" & "10°18'25.78""E" \\ \hline
        Corner Upper Left Lat DMS & "37°05'44.84""N" & "37°05'44.84""N" & "37°05'44.88""N" \\ \hline
        Corner Upper Left Long DMS & "9°03'06.34""E" & "9°02'54.20""E" & "9°02'05.57""E" \\ \hline
        Corner Upper Right Lat DMS & "37°03'59.90""N" & "37°04'00.44""N" & "37°04'01.24""N" \\ \hline
        Corner Upper Right Long DMS & "11°37'46.99""E" & "11°37'22.69""E" & "11°36'46.30""E" \\ \hline
        Corner Lower Left Lat DMS & "34°59'48.84""N" & "34°59'48.84""N" & "34°59'48.84""N" \\ \hline
        Corner Lower Left Long DMS & "9°03'01.48""E" & "9°02'49.63""E" & "9°02'02.29""E" \\ \hline
        Corner Lower Right Lat DMS & "34°58'11.60""N" & "34°58'12.11""N" & "34°58'12.86""N" \\ \hline
        Corner Lower Right Long DMS & "11°33'39.28""E" & "11°33'15.66""E" & "11°32'40.20""E" \\ \hline
        Scene Center Latitude & 36.04314 & 36.04311 & 36.04327 \\ \hline
        Scene Center Longitude & 10.32396 & 10.31843 & 10.30716 \\ \hline
        Corner Upper Left Latitude & 37.09579 & 37.09579 & 37.09580 \\ \hline
        Corner Upper Left Longitude & 9.05176 & 9.04839 & 9.03488 \\ \hline
        Corner Upper Right Latitude & 37.06664 & 37.06679 & 37.06701 \\ \hline
        Corner Upper Right Longitude & 11.62972 & 11.62297 & 11.61286 \\ \hline
        Corner Lower Left Latitude & 34.99690 & 34.99690 & 34.99690 \\ \hline
        Corner Lower Left Longitude & 9.05041 & 9.04712 & 9.03397 \\ \hline
        Corner Lower Right Latitude & 34.96989 & 34.97003 & 34.97024 \\ \hline
        Corner Lower Right Longitude & 11.56091 & 11.55435 & 11.54450 \\ 
        \bottomrule
 \end{tabularx}
 \end{adjustwidth}
\end{table}

\begin{table}[H]
\scriptsize
\caption{Metadata of Landsat OLI/TIRS satellite images for the Gulf of Hammamet area (year 2017).\label{tabApp}}
\begin{adjustwidth}{-\extralength}{0.5cm}
\centering
\begin{tabularx}{\fulllength}{|l|l|l|l}
    \toprule
      Data Set Attribute & Attribute Value & Attribute Value & Attribute Value \\ \hline
        Landsat Scene Identifier & LC81910352017187LGN00 & LC81910352017091LGN00 & LC81910352017059LGN00 \\ \hline
        Date Acquired & 2017/07/06 & 2017/04/01 & 2017/02/28 \\ \hline
        Collection Category & T1 & T1 & T1 \\ \hline
        Collection Number & 2 & 2 & 2 \\ \hline
        WRS Path & 191 & 191 & 191 \\ \hline
        WRS Row & 35 & 35 & 35 \\ \hline
        Target WRS Path & 191 & 191 & 191 \\ \hline
        Target WRS Row & 35 & 35 & 35 \\ \hline
        Nadir/Off Nadir & NADIR & NADIR & NADIR \\ \hline
        Roll Angle & 0.000 & 0.000 & 0.000 \\ \hline
        Date Product Generated L2 & 2020/09/03 & 2020/09/04 & 2020/09/05 \\ \hline
        Date Product Generated L1 & 2020/09/03 & 2020/09/04 & 2020/09/05 \\ \hline
        Start Time & 2017-07-06 09:54:30.064594 & 2017-04-01 09:54:16.946987 & 2017-02-28 09:54:34.167045 \\ \hline
        Stop Time & 2017-07-06 09:55:01.834592 & 2017-04-01 09:54:48.716984 & 2017-02-28 09:55:05.937042 \\ \hline
        Station Identifier & LGN & LGN & LGN \\ \hline
        Day/Night Indicator & DAY & DAY & DAY \\ \hline
        Land Cloud Cover & 0.02 & 0.49 & 0.58 \\ \hline
        Scene Cloud Cover L1 & 3.34 & 0.36 & 0.45 \\ \hline
        Ground Control Points Model & 1168 & 1112 & 1111 \\ \hline
        Ground Control Points Version & 5 & 5 & 5 \\ \hline
        Geometric RMSE Model & 5.176 & 5.628 & 5.933 \\ \hline
        Geometric RMSE Model X & 3.614 & 3.874 & 4.125 \\ \hline
        Geometric RMSE Model Y & 3.705 & 4.082 & 4.264 \\ \hline
        Processing Software Version & LPGS\_15.3.1c & LPGS\_15.3.1c & LPGS\_15.3.1c \\ \hline
        Sun Elevation L0RA & 66.59843288 & 52.68839245 & 40.64086896 \\ \hline
        Sun Azimuth L0RA & 118.86194381 & 142.00904413 & 147.81922946 \\ \hline
        TIRS SSM Model & FINAL & FINAL & FINAL \\ \hline
        Data Type L2 & OLI\_TIRS\_L2SP & OLI\_TIRS\_L2SP & OLI\_TIRS\_L2SP \\ \hline
        Sensor Identifier & OLI\_TIRS & OLI\_TIRS & OLI\_TIRS \\ \hline
        Satellite & 8 & 8 & 8 \\ \hline
        Product Map Projection L1 & UTM & UTM & UTM \\ \hline
        UTM Zone & 32 & 32 & 32 \\ \hline
        Datum & WGS84 & WGS84 & WGS84 \\ \hline
        Ellipsoid & WGS84 & WGS84 & WGS84 \\ \hline
        Scene Center Lat DMS & "36°02'36.53""N" & "36°02'36.49""N" & "36°02'35.63""N" \\ \hline
        Scene Center Long DMS & "10°21'44.68""E" & "10°20'01.54""E" & "10°20'13.45""E" \\ \hline
        Corner Upper Left Lat DMS & "37°05'54.49""N" & "37°05'44.84""N" & "37°05'44.81""N" \\ \hline
        Corner Upper Left Long DMS & "9°05'32.21""E" & "9°03'42.80""E" & "9°03'54.97""E" \\ \hline
        Corner Upper Right Lat DMS & "37°04'06.64""N" & "37°03'59.08""N" & "37°03'58.82""N" \\ \hline
        Corner Upper Right Long DMS & "11°40'00.80""E" & "11°38'23.39""E" & "11°38'35.52""E" \\ \hline
        Corner Lower Left Lat DMS & "34°59'48.73""N" & "34°59'39.08""N" & "34°59'39.05""N" \\ \hline
        Corner Lower Left Long DMS & "9°05'23.50""E" & "9°03'36.97""E" & "9°03'48.82""E" \\ \hline
        Corner Lower Right Lat DMS & "34°58'08.83""N" & "34°58'01.13""N" & "34°58'00.88""N" \\ \hline
        Corner Lower Right Long DMS & "11°35'49.31""E" & "11°34'14.45""E" & "11°34'26.26""E" \\ \hline
        Scene Center Latitude & 36.04348 & 36.04347 & 36.04323 \\ \hline
        Scene Center Longitude & 10.36241 & 10.33376 & 10.33707 \\ \hline
        Corner Upper Left Latitude & 37.09847 & 37.09579 & 37.09578 \\ \hline
        Corner Upper Left Longitude & 9.09228 & 9.06189 & 9.06527 \\ \hline
        Corner Upper Right Latitude & 37.06851 & 37.06641 & 37.06634 \\ \hline
        Corner Upper Right Longitude & 11.66689 & 11.63983 & 11.64320 \\ \hline
        Corner Lower Left Latitude & 34.99687 & 34.99419 & 34.99418 \\ \hline
        Corner Lower Left Longitude & 9.08986 & 9.06027 & 9.06356 \\ \hline
        Corner Lower Right Latitude & 34.96912 & 34.96698 & 34.96691 \\ \hline
        Corner Lower Right Longitude & 11.59703 & 11.57068 & 11.57396 \\ 
        \bottomrule 
 \end{tabularx}
 \end{adjustwidth}
\end{table}

\subsection[\appendixname~\thesubsection]{Gulf of Gabès area}

\begin{table}[H]
\scriptsize
\caption{Metadata of Landsat OLI/TIRS satellite images for the Gulf of Gabès area (year 2023).\label{tabApp}}
\begin{adjustwidth}{-\extralength}{0.5cm}
\centering
\begin{tabularx}{\fulllength}{|l|l|l|l}
    \toprule
      Data Set Attribute & Attribute Value & Attribute Value & Attribute Value \\ \hline
        Landsat Scene Identifier & LC91910362023196LGN00 & LC91910362023100LGN00 & LC91910362023052LGN01 \\ \hline
        Date Acquired & 2023/07/15 & 2023/04/10 & 2023/02/21 \\ \hline
        Collection Category & T1 & T1 & T1 \\ \hline
        Collection Number & 2 & 2 & 2 \\ \hline
        WRS Path & 191 & 191 & 191 \\ \hline
        WRS Row & 36 & 36 & 36 \\ \hline
        Target WRS Path & 191 & 191 & 191 \\ \hline
        Target WRS Row & 36 & 36 & 36 \\ \hline
        Nadir/Off Nadir & NADIR & NADIR & NADIR \\ \hline
        Roll Angle & 0.001 & 0.000 & -0.001 \\ \hline
        Date Product Generated L2 & 2023/07/17 & 2023/04/12 & 2023/03/09 \\ \hline
        Date Product Generated L1 & 2023/07/15 & 2023/04/10 & 2023/03/09 \\ \hline
        Start Time & 2023-07-15 09:54:42 & 2023-04-10 09:55:00 & 2023-02-21 09:55:26 \\ \hline
        Stop Time & 2023-07-15 09:55:14 & 2023-04-10 09:55:32 & 2023-02-21 09:55:57 \\ \hline
        Station Identifier & LGN & LGN & LGN \\ \hline
        Day/Night Indicator & DAY & DAY & DAY \\ \hline
        Land Cloud Cover & 0.00 & 0.14 & 0.11 \\ \hline
        Scene Cloud Cover L1 & 1.04 & 0.17 & 0.13 \\ \hline
        Ground Control Points Model & 1131 & 1140 & 1147 \\ \hline
        Ground Control Points Version & 5 & 5 & 5 \\ \hline
        Geometric RMSE Model & 5.456 & 5.149 & 5.735 \\ \hline
        Geometric RMSE Model X & 3.725 & 3.659 & 4.178 \\ \hline
        Geometric RMSE Model Y & 3.986 & 3.624 & 3.929 \\ \hline
        Processing Software Version & LPGS\_16.3.0 & LPGS\_16.2.0 & LPGS\_16.2.0 \\ \hline
        Sun Elevation L0RA & 66.06839265 & 56.73154182 & 39.15035895 \\ \hline
        Sun Azimuth L0RA & 117.11855615 & 138.30347738 & 148.18246509 \\ \hline
        TIRS SSM Model & N/A & N/A & N/A \\ \hline
        Data Type L2 & OLI\_TIRS\_L2SP & OLI\_TIRS\_L2SP & OLI\_TIRS\_L2SP \\ \hline
        Sensor Identifier & OLI\_TIRS & OLI\_TIRS & OLI\_TIRS \\ \hline
        Satellite & 9 & 9 & 9 \\ \hline
        Product Map Projection L1 & UTM & UTM & UTM \\ \hline
        UTM Zone & 32 & 32 & 32 \\ \hline
        Datum & WGS84 & WGS84 & WGS84 \\ \hline
        Ellipsoid & WGS84 & WGS84 & WGS84 \\ \hline
        Scene Center Lat DMS & "34°36'39.31""N" & "34°36'39.49""N" & "34°36'38.20""N" \\ \hline
        Scene Center Long DMS & "9°54'48.85""E" & "9°54'34.67""E" & "9°53'51.79""E" \\ \hline
        Corner Upper Left Lat DMS & "35°39'42.73""N" & "35°39'32.94""N" & "35°39'32.83""N" \\ \hline
        Corner Upper Left Long DMS & "8°39'23.04""E" & "8°39'11.12""E" & "8°38'35.34""E" \\ \hline
        Corner Upper Right Lat DMS & "35°38'32.71""N" & "35°38'23.21""N" & "35°38'24.07""N" \\ \hline
        Corner Upper Right Long DMS & "11°11'28.68""E" & "11°11'16.48""E" & "11°10'28.81""E" \\ \hline
        Corner Lower Left Lat DMS & "33°33'25.56""N" & "33°33'25.52""N" & "33°33'25.42""N" \\ \hline
        Corner Lower Left Long DMS & "8°39'53.86""E" & "8°39'42.23""E" & "8°39'07.31""E" \\ \hline
        Corner Lower Right Lat DMS & "33°32'20.80""N" & "33°32'21.01""N" & "33°32'21.80""N" \\ \hline
        Corner Lower Right Long DMS & "11°08'12.23""E" & "11°08'00.60""E" & "11°07'14.12""E" \\ \hline
        Scene Center Latitude & 34.61092 & 34.61097 & 34.61061 \\ \hline
        Scene Center Longitude & 9.91357 & 9.90963 & 9.89772 \\ \hline
        Corner Upper Left Latitude & 35.66187 & 35.65915 & 35.65912 \\ \hline
        Corner Upper Left Longitude & 8.65640 & 8.65309 & 8.64315 \\ \hline
        Corner Upper Right Latitude & 35.64242 & 35.63978 & 35.64002 \\ \hline
        Corner Upper Right Longitude & 11.19130 & 11.18791 & 11.17467 \\ \hline
        Corner Lower Left Latitude & 33.55710 & 33.55709 & 33.55706 \\ \hline
        Corner Lower Left Longitude & 8.66496 & 8.66173 & 8.65203 \\ \hline
        Corner Lower Right Latitude & 33.53911 & 33.53917 & 33.53939 \\ \hline
        Corner Lower Right Longitude & 11.13673 & 11.13350 & 11.12059 \\ \hline
     \bottomrule
 \end{tabularx}
 \end{adjustwidth}
\end{table}

\begin{table}[H]
\scriptsize
\caption{Metadata of Landsat OLI/TIRS satellite images for the Gulf of Gabès area (year 2017).\label{tabApp}}
\begin{adjustwidth}{-\extralength}{0.5cm}
\centering
\begin{tabularx}{\fulllength}{|l|l|l|l}
    \toprule
     Data Set Attribute & Attribute Value & Attribute Value & Attribute Value \\ \hline
        Landsat Scene Identifier & LC81910362017203LGN00 & LC81910362017091LGN00 & LC81910362017059LGN00 \\ \hline
        Date Acquired & 2017/07/22 & 2017/04/01 & 2017/02/28 \\ \hline
        Collection Category & T1 & T1 & T1 \\ \hline
        Collection Number & 2 & 2 & 2 \\ \hline
        WRS Path & 191 & 191 & 191 \\ \hline
        WRS Row & 36 & 36 & 36 \\ \hline
        Target WRS Path & 191 & 191 & 191 \\ \hline
        Target WRS Row & 36 & 36 & 36 \\ \hline
        Nadir/Off Nadir & NADIR & NADIR & NADIR \\ \hline
        Roll Angle & 0.000 & 0.000 & 0.000 \\ \hline
        Date Product Generated L2 & 2020/09/03 & 2020/09/04 & 2020/09/05 \\ \hline
        Date Product Generated L1 & 2020/09/03 & 2020/09/04 & 2020/09/05 \\ \hline
        Start Time & 2017-07-22 09:54:59.645678 & 2017-04-01 09:54:40.842261 & 2017-02-28 09:54:58.053847 \\ \hline
        Stop Time & 2017-07-22 09:55:31.415676 & 2017-04-01 09:55:12.612258 & 2017-02-28 09:55:29.823844 \\ \hline
        Station Identifier & LGN & LGN & LGN \\ \hline
        Day/Night Indicator & DAY & DAY & DAY \\ \hline
        Land Cloud Cover & 2.48 & 1.62 & 0.00 \\ \hline
        Scene Cloud Cover L1 & 3.07 & 1.39 & 0.00 \\ \hline
        Ground Control Points Model & 1108 & 1146 & 1153 \\ \hline
        Ground Control Points Version & 5 & 5 & 5 \\ \hline
        Geometric RMSE Model & 4.771 & 4.844 & 5.169 \\ \hline
        Geometric RMSE Model X & 3.282 & 3.393 & 3.675 \\ \hline
        Geometric RMSE Model Y & 3.462 & 3.458 & 3.635 \\ \hline
        Processing Software Version & LPGS\_15.3.1c & LPGS\_15.3.1c & LPGS\_15.3.1c \\ \hline
        Sun Elevation L0RA & 65.21568136 & 53.64685927 & 41.70987822 \\ \hline
        Sun Azimuth L0RA & 119.62146872 & 140.36767543 & 146.77696078 \\ \hline
        TIRS SSM Model & FINAL & FINAL & FINAL \\ \hline
        Data Type L2 & OLI\_TIRS\_L2SP & OLI\_TIRS\_L2SP & OLI\_TIRS\_L2SP \\ \hline
        Sensor Identifier & OLI\_TIRS & OLI\_TIRS & OLI\_TIRS \\ \hline
        Satellite & 8 & 8 & 8 \\ \hline
        Product Map Projection L1 & UTM & UTM & UTM \\ \hline
        UTM Zone & 32 & 32 & 32 \\ \hline
        Datum & WGS84 & WGS84 & WGS84 \\ \hline
        Ellipsoid & WGS84 & WGS84 & WGS84 \\ \hline
        Scene Center Lat DMS & "34°36'37.94""N" & "34°36'37.76""N" & "34°36'38.84""N" \\ \hline
        Scene Center Long DMS & "9°57'06.12""E" & "9°55'28.24""E" & "9°55'39.86""E" \\ \hline
        Corner Upper Left Lat DMS & "35°39'43.09""N" & "35°39'42.84""N" & "35°39'42.88""N" \\ \hline
        Corner Upper Left Long DMS & "8°41'46.21""E" & "8°40'10.74""E" & "8°40'22.69""E" \\ \hline
        Corner Upper Right Lat DMS & "35°38'30.30""N" & "35°38'32.06""N" & "35°38'31.85""N" \\ \hline
        Corner Upper Right Long DMS & "11°13'39.79""E" & "11°12'04.43""E" & "11°12'16.34""E" \\ \hline
        Corner Lower Left Lat DMS & "33°33'25.92""N" & "33°33'25.67""N" & "33°33'25.70""N" \\ \hline
        Corner Lower Left Long DMS & "8°42'13.46""E" & "8°40'40.40""E" & "8°40'52.03""E" \\ \hline
        Corner Lower Right Lat DMS & "33°32'18.56""N" & "33°32'20.18""N" & "33°32'20""N" \\ \hline
        Corner Lower Right Long DMS & "11°10'20.10""E" & "11°08'47.11""E" & "11°08'58.74""E" \\ \hline
        Scene Center Latitude & 34.61054 & 34.61049 & 34.61079 \\ \hline
        Scene Center Longitude & 9.95170 & 9.92451 & 9.92774 \\ \hline
        Corner Upper Left Latitude & 35.66197 & 35.66190 & 35.66191 \\ \hline
        Corner Upper Left Longitude & 8.69617 & 8.66965 & 8.67297 \\ \hline
        Corner Upper Right Latitude & 35.64175 & 35.64224 & 35.64218 \\ \hline
        Corner Upper Right Longitude & 11.22772 & 11.20123 & 11.20454 \\ \hline
        Corner Lower Left Latitude & 33.55720 & 33.55713 & 33.55714 \\ \hline
        Corner Lower Left Longitude & 8.70374 & 8.67789 & 8.68112 \\ \hline
        Corner Lower Right Latitude & 33.53849 & 33.53894 & 33.53889 \\ \hline
        Corner Lower Right Longitude & 11.17225 & 11.14642 & 11.14965 \\ \hline 
     \bottomrule
 \end{tabularx}
 \end{adjustwidth}
\end{table}

\pagebreak
\section[\appendixname~\thesection]{Clustering reports}\label{AppB}

\subsection[\appendixname~\thesubsection]{Gulf of Hammamet area}

\subsubsection[\appendixname~\thesubsubsection]{February 2017}

\begin{scriptsize}
\begin{verbatim}
#################### CLUSTER (Thu Jul 20 19:42:28 2023) ####################
Location: Tunisia
Mapset:   PERMANENT
Group:    L8_2017
Subgroup: res_30m
 L8_2017_01@PERMANENT
 L8_2017_02@PERMANENT
 L8_2017_03@PERMANENT
 L8_2017_04@PERMANENT
 L8_2017_05@PERMANENT
 L8_2017_06@PERMANENT
 L8_2017_07@PERMANENT
Result signature file: cluster_L8_2017

Region
  North:   4105515.00  East:    735015.00
  South:   3872385.00  West:    505785.00
  Res:          30.00  Res:         30.00
  Rows:          7771  Cols:         7641  Cells: 59378211
Mask: no
Cluster parameters
  Nombre de classes initiales: 	10
 Minimum class size:           17
 Minimum class separation:     0.000000
 Percent convergence:          98.000000
 Maximum number of iterations: 30
 Row sampling interval:        77
 Col sampling interval:        76
Sample size: 6933 points
means and standard deviations for 7 bands
moyennes 8708.15 9123.24 10478.6 10923.3 16084 14664.6 12529.6
écart-type 1111.56 1367.1 2328.79 3138.4 5990.96 5443.95 4286.03
initial means for each band

classe 1    7596.58 7756.14 8149.85 7784.86 10093 9220.67 8243.53
classe 2    7843.6 8059.94 8667.35 8482.28 11424.3 10430.4 9195.98
classe 3    8090.61 8363.74 9184.86 9179.7 12755.7 11640.2 10148.4
classe 4    8337.63 8667.54 9702.37 9877.13 14087 12850 11100.9
classe 5    8584.64 8971.34 10219.9 10574.5 15418.3 14059.7 12053.3
classe 6    8831.65 9275.14 10737.4 11272 16749.6 15269.5 13005.8
classe 7    9078.67 9578.94 11254.9 11969.4 18081 16479.3 13958.2
classe 8    9325.68 9882.74 11772.4 12666.8 19412.3 17689 14910.7
classe 9    9572.7 10186.5 12289.9 13364.2 20743.6 18898.8 15863.1
classe 10   9819.71 10490.3 12807.4 14061.7 22074.9 20108.6 16815.6

class means/stddev for each band
class 1 (1936)
moyennes 7543.95 7678.06 7642.34 7372.52 7339.36 7470.75 7460.05
écart-type 280.142 343.706 664.799 581.133 305.826 183.653 145.669
class 2 (63)
moyennes 9034.81 9522.65 11242.6 11436.8 10405.2 8682.76 8118.87
écart-type 1129.8 1375.17 2395.17 2825.66 1627.25 1370.64 803.186
class 3 (88)
moyennes 8222 8512.48 9550.11 9651.31 13695.8 10865.3 9397.86
écart-type 633.751 818.332 1556.44 2098.95 1065.63 1491.64 962.104
class 4 (159)
moyennes 8325.57 8639.04 9760.28 9824.18 15274.9 12332.6 10315.3
écart-type 583.028 753.624 1348.79 1813.64 1684.3 1428.11 1135.42
class 5 (375)
moyennes 8234.31 8505.45 9671.42 9482.12 18154.5 13424.7 10856.3
écart-type 410.004 479.463 611.068 1030.31 2647.22 1090.87 1202.22
class 6 (854)
moyennes 8387.59 8714.15 10142.5 10017.2 19626.9 14529.1 11619.5
écart-type 440.904 530.035 665.934 1233.49 2971.53 1303.13 1507.6
class 7 (914)
moyennes 8751.68 9171.15 10830.1 11062.4 19784.6 16091.2 13036.9
écart-type 530.04 623.566 780.397 1379.93 3179.75 1377.44 1672.65
class 8 (750)
moyennes 9180.8 9728.57 11647.7 12414.3 19413.2 17923.2 14861.3
écart-type 504.969 580.175 697.874 1069.45 2273.84 1087.71 1344.09
class 9 (642)
moyennes 9567.56 10208.8 12368.6 13571.8 19616.8 19489.8 16393.6
écart-type 508.989 595.323 655.666 894.536 1676.44 833.279 1100.4
class 10 (1152)
moyennes 10307.4 11110.9 13792.2 15694.9 21011.7 22255.3 18980.5
écart-type 916.372 1065.89 1311.89 1549.92 1621.65 1735.38 1767.02

Distribution des classes
       1936         63         88        159        375
        854        914        750        642       1152
######## iteration 1 ###########
10 classes, 73.11% points stable
Distribution des classes
       1896        105         94        260        278
        855        687        904        969        885
######## iteration 2 ###########
10 classes, 88.89% points stable
Distribution des classes
       1901         98        121        334        273
        744        719        917       1080        746
######## iteration 3 ###########
10 classes, 93.81% points stable
Distribution des classes
       1904         95        148        378        354
        624        726        890       1143        671
######## iteration 4 ###########
10 classes, 95.28% points stable
Distribution des classes
       1906         94        177        418        413
        546        717        872       1175        615
######## iteration 5 ###########
10 classes, 96.37% points stable
Distribution des classes
       1907         95        193        476        466
        493        684        855       1184        580
######## iteration 6 ###########
10 classes, 96.68% points stable
Distribution des classes
       1909         96        214        522        508
        457        642        862       1169        554
######## iteration 7 ###########
10 classes, 97.35% points stable
Distribution des classes
       1910         96        234        567        529
        439        605        874       1145        534
######## iteration 8 ###########
10 classes, 97.72% points stable
Distribution des classes
       1910         96        252        605        543
        428        583        875       1131        510
######## iteration 9 ###########
10 classes, 98.02% points stable
Distribution des classes
       1911         96        267        631        557
        424        562        884       1108        493
########## final results #############
10 classes (convergence=98.0%)

class separability matrix
         1     2     3     4     5     6     7     8     9    10

  1      0
  2    2.6     0
  3    3.0   2.2     0
  4    4.9   2.6   1.3     0
  5    5.0   2.6   1.0   1.1     0
  6    4.6   2.6   1.3   1.7   0.7     0
  7    6.0   3.0   2.0   0.8   1.4   1.8     0
  8    6.5   3.4   2.3   1.0   2.0   2.6   0.8     0
  9    7.6   4.0   3.3   1.9   3.0   3.4   1.6   1.0     0
 10    7.0   4.3   3.7   2.6   3.4   3.7   2.2   1.8   1.0     0


class means/stddev for each band
class 1 (1911)
moyennes 7523.69 7650.46 7586.15 7317.79 7326.28 7470.52 7460.71
écart-type 199.208 226.738 433.861 315.139 276.934 192.008 154.064
class 2 (96)
moyennes 9603.39 10327.9 12814 13338.7 10794.6 7933.64 7617.12
écart-type 673.173 704.094 1079.53 1652.58 2943.71 971.297 514.661
class 3 (267)
moyennes 8108.04 8348.92 9190.7 9132.8 14466 12328.1 10382.7
écart-type 279.812 320.437 518.45 663.449 1455.6 1163.5 923.442
class 4 (631)
moyennes 8751.79 9156.74 10603.9 11090.7 16711.6 15600.1 12927.6
écart-type 395.654 463.445 709.094 993.703 1279.39 913.201 974.364
class 5 (557)
moyennes 8182.17 8475.2 9896.99 9464.78 20356.9 13922.5 10886.8
écart-type 260.723 292.6 465.421 623.226 1188.81 1021.22 808.439
class 6 (424)
moyennes 8030.23 8251.04 9628.57 8779.63 24296.8 13239.5 10064.6
écart-type 267.67 282.628 501.637 592.081 1692.55 951.726 713.051
class 7 (562)
moyennes 8881.47 9383.44 11280.9 11477.3 21214.6 17035.3 13592.8
écart-type 359.019 407.881 554.181 778.486 1512.62 1030.69 919.513
class 8 (884)
moyennes 9275.02 9820.84 11668.6 12711.5 17849.5 18128.5 15313.9
écart-type 454.737 536.129 732.941 843.276 1234.67 913.051 957.389
class 9 (1108)
moyennes 9833.81 10533.4 12878.9 14434.5 19815.8 20584.5 17484.2
écart-type 512.17 596.387 700.693 857.892 1412.56 1004.99 1082.09
class 10 (493)
moyennes 10827.6 11739.6 14718.5 16945.9 21841.6 23775.4 20423.6
écart-type 1094.62 1257.14 1425.55 1419.46 1500.15 1449.78 1461.87

#################### CLASSES ####################
10 classes, 98.02% points stable

######## CLUSTER END (Thu Jul 20 19:42:28 2023) ########
\end{verbatim}
\end{scriptsize}

\subsubsection[\appendixname~\thesubsubsection]{April 2017}
\begin{scriptsize}
\begin{verbatim}
#################### CLUSTER (Thu Jul 20 20:20:27 2023) ####################
Location: Tunisia
Mapset:   PERMANENT
Group:    L8_2017a
Subgroup: res_30m
 L8_2017a_01@PERMANENT
 L8_2017a_02@PERMANENT
 L8_2017a_03@PERMANENT
 L8_2017a_04@PERMANENT
 L8_2017a_05@PERMANENT
 L8_2017a_06@PERMANENT
 L8_2017a_07@PERMANENT
Result signature file: cluster_L8_2017a
Region
  North:   4105515.00  East:    735015.00
  South:   3872385.00  West:    505785.00
  Res:          30.00  Res:         30.00
  Rows:          7771  Cols:         7641  Cells: 59378211
Mask: no
Cluster parameters
  Nombre de classes initiales: 	10
 Minimum class size:           17
 Minimum class separation:     0.000000
 Percent convergence:          98.000000
 Maximum number of iterations: 30
 Row sampling interval:        77
 Col sampling interval:        76

Sample size: 6931 points

means and standard deviations for 7 bands

moyennes 8910.31 9388.99 10819.6 11430.1 16494.7 15518.8 13323.3
écart-type 1277.01 1555.04 2548.9 3403.47 6032.48 5822.77 4672.21

initial means for each band

classe 1    7633.3 7833.95 8270.68 8026.63 10462.2 9696.02 8651.08
classe 2    7917.08 8179.51 8837.1 8782.96 11802.8 10990 9689.35
classe 3    8200.86 8525.08 9403.52 9539.28 13143.3 12283.9 10727.6
classe 4    8484.64 8870.64 9969.94 10295.6 14483.9 13577.9 11765.9
classe 5    8768.42 9216.21 10536.4 11051.9 15824.4 14871.8 12804.2
classe 6    9052.2 9561.77 11102.8 11808.3 17165 16165.8 13842.4
classe 7    9335.98 9907.34 11669.2 12564.6 18505.5 17459.7 14880.7
classe 8    9619.75 10252.9 12235.6 13320.9 19846.1 18753.7 15919
classe 9    9903.53 10598.5 12802.1 14077.2 21186.6 20047.6 16957.2
classe 10   10187.3 10944 13368.5 14833.6 22527.2 21341.6 17995.5
class means/stddev for each band
class 1 (1936)
moyennes 7574.24 7735.09 7721.36 7466.68 7414.89 7628.6 7630.8
écart-type 367.279 449.271 852.879 842.89 305.122 210.854 192.701
class 2 (28)
moyennes 8971.82 9485.46 11084.3 11131.1 11548.5 9228.86 8476.14
écart-type 1480.91 1647.17 2344.9 2529.6 1904.7 1480.2 838.668
class 3 (75)
moyennes 8406.47 8682.52 9622.69 9557.69 14226 11871.1 10019.7
écart-type 1173.11 1311.78 1785.02 2011.23 1324.86 1220.47 784.185
class 4 (132)
moyennes 8287.05 8560.14 9566.78 9489.11 16401.2 13403.8 11004.4
écart-type 299.277 339.017 477.664 641.746 1860.08 789.263 840.708
class 5 (528)
moyennes 8299.29 8603.93 9868.33 9691.18 19692.1 14034 11154.8
écart-type 513.459 571.896 722.457 1061.18 2732.1 1132.36 1211.84
class 6 (845)
moyennes 8617.56 9017.26 10536.6 10630.1 20110.4 15444.9 12322.9
écart-type 600.136 689.469 870.244 1337.8 2747.94 1267.86 1445.68
class 7 (808)
moyennes 9155.82 9680.96 11462.2 12146.5 19283 17521.7 14418.2
écart-type 801.789 860.707 982.654 1273.23 2093.61 1274.61 1430.06
class 8 (751)
moyennes 9509.78 10144.9 12161.5 13215.5 19632.3 19036 15950.4
écart-type 763.187 829.15 920.492 1069.43 1635.46 1080 1234.81
class 9 (614)
moyennes 9833.67 10546.6 12796.8 14192.9 20152.5 20570.3 17491.8
écart-type 653.451 725.231 762.044 830.39 1338.62 865.407 1083.79
class 10 (1214)
moyennes 10606.7 11510.8 14317.3 16418.8 21684.7 23335.2 20146.5
écart-type 941.57 1108.14 1325.98 1480.89 1481.03 1619.68 1676.31

Distribution des classes
       1936         28         75        132        528
        845        808        751        614       1214
######## iteration 1 ###########
10 classes, 77.30% points stable
Distribution des classes
       1890         75         76        245        613
        531        897        759        888        957
######## iteration 2 ###########
10 classes, 91.52% points stable
Distribution des classes
       1895         75        102        317        577
        543        781        869        958        814
######## iteration 3 ###########
10 classes, 94.94% points stable
Distribution des classes
       1897         81        127        335        565
        571        721        941        972        721
######## iteration 4 ###########
10 classes, 96.25% points stable
Distribution des classes
       1898         80        141        354        573
        576        705        968        984        652
######## iteration 5 ###########
10 classes, 97.13% points stable
Distribution des classes
       1898         79        148        368        588
        581        693        981        998        597
######## iteration 6 ###########
10 classes, 97.50% points stable
Distribution des classes
       1899         78        155        381        598
        586        696        980        997        561
######## iteration 7 ###########
10 classes, 97.52% points stable
Distribution des classes
       1900         78        157        398        612
        581        698        998        980        529
######## iteration 8 ###########
10 classes, 97.73% points stable
Distribution des classes
       1900         79        162        409        625
        577        708       1005        967        499
######## iteration 9 ###########
10 classes, 97.82% points stable
Distribution des classes
       1900         79        172        424        635
        573        712       1004        956        476
######## iteration 10 ###########
10 classes, 98.12% points stable
Distribution des classes
       1901         80        178        440        641
        571        715       1006        939        460
########## final results #############
10 classes (convergence=98.1%)
class separability matrix

         1     2     3     4     5     6     7     8     9    10

  1      0
  2    1.7     0
  3    3.1   1.7     0
  4    5.3   2.2   1.1     0
  5    4.7   1.9   1.0   1.3     0
  6    6.3   2.4   1.8   0.9   1.6     0
  7    5.9   2.4   1.9   1.1   2.2   0.8     0
  8    7.6   3.0   2.9   2.1   3.1   1.4   0.8     0
  9    8.5   3.4   3.8   3.0   4.0   2.3   1.7   1.0     0
 10    7.9   3.7   4.2   3.5   4.2   2.9   2.4   1.8   1.0     0

class means/stddev for each band
class 1 (1901)
moyennes 7537.46 7687.73 7626.91 7363.61 7399.07 7630.31 7633.45
écart-type 243.321 281.004 496.139 363.319 278.956 213.743 194.02
class 2 (80)
moyennes 11384.7 12171.8 14674.2 15354.2 12589.9 8667.99 8091.54
écart-type 2668.9 2685.69 2569.09 3106.74 5398.79 1721.66 851.413
class 3 (178)
moyennes 8218.57 8481.5 9398.38 9325.25 14998.5 12698 10601.1
écart-type 294.276 340.251 575.419 677.483 1362.72 1168.29 893.63
class 4 (440)
moyennes 8633.37 9011.06 10342.6 10592.5 17502.2 15377 12581.1
écart-type 291.159 339.709 513.635 677.724 1361.55 767.388 781.177
class 5 (641)
moyennes 8116.34 8405.74 9757.32 9277.91 22418.5 13760.1 10568.5
écart-type 220.641 241.767 431.078 571.397 1740.46 842.148 652.567
class 6 (571)
moyennes 8829.79 9337.73 11140.6 11376.9 21187.7 16890.9 13374
écart-type 326.563 375.923 522.283 705.858 1184.33 993.323 885.803
class 7 (715)
moyennes 9227.36 9751.14 11468.8 12381.4 17758.6 17812.1 14959.4
écart-type 499.514 588.499 782.755 954.599 1197.54 920.169 1060.86
class 8 (1006)
moyennes 9615.55 10281.3 12393.7 13625.2 19660.5 19779.2 16697.3
écart-type 466.391 539.252 667.241 735.077 1418.28 826.435 960.02
class 9 (939)
moyennes 10211.7 11025.3 13578.6 15438.8 20743.3 22114.8 19060.1
écart-type 562.004 668.667 769.605 788.966 1264.55 957.686 1063.92
class 10 (460)
moyennes 11200.5 12227.5 15347.5 17771.3 22736.8 24938.5 21653.6
écart-type 1156.3 1338.64 1449.63 1393.37 1328.07 1276.32 1383.58
#################### CLASSES ####################
10 classes, 98.12% points stable

######## CLUSTER END (Thu Jul 20 20:20:27 2023) ########
\end{verbatim}
\end{scriptsize}

\subsubsection[\appendixname~\thesubsubsection]{July 2017}

\begin{scriptsize}
\begin{verbatim}
#################### CLUSTER (Thu Jul 20 20:31:04 2023) ####################
Location: Tunisia
Mapset:   PERMANENT
Group:    L8_2017j
Subgroup: res_30m
 L8_2017j_01@PERMANENT
 L8_2017j_02@PERMANENT
 L8_2017j_03@PERMANENT
 L8_2017j_04@PERMANENT
 L8_2017j_05@PERMANENT
 L8_2017j_06@PERMANENT
 L8_2017j_07@PERMANENT
Result signature file: cluster_L8_2017j

Region
  North:   4105515.00  East:    735015.00
  South:   3872685.00  West:    508155.00
  Res:          30.00  Res:         30.00
  Rows:          7761  Cols:         7562  Cells: 58688682
Mask: no

Cluster parameters
  Nombre de classes initiales: 	10
 Minimum class size:           17
 Minimum class separation:     0.000000
 Percent convergence:          98.000000
 Maximum number of iterations: 30
 Row sampling interval:        77
 Col sampling interval:        75
Sample size: 7026 points

means and standard deviations for 7 bands

moyennes 9904.03 10666.9 12629.3 14200.5 17375.7 19128.8 16470.7
écart-type 1966.49 2202.44 3155.47 4346.23 5675.2 6751.68 5230.25

initial means for each band
classe 1    7937.54 8464.42 9473.87 9854.23 11700.5 12377.1 11240.4
classe 2    8374.54 8953.85 10175.1 10820.1 12961.6 13877.5 12402.7
classe 3    8811.54 9443.28 10876.3 11785.9 14222.8 15377.9 13565
classe 4    9248.54 9932.71 11577.5 12751.7 15483.9 16878.2 14727.3
classe 5    9685.53 10422.1 12278.7 13717.5 16745.1 18378.6 15889.5
classe 6    10122.5 10911.6 12979.9 14683.4 18006.2 19879 17051.8
classe 7    10559.5 11401 13681.2 15649.2 19267.4 21379.4 18214.1
classe 8    10996.5 11890.4 14382.4 16615 20528.6 22879.7 19376.4
classe 9    11433.5 12379.9 15083.6 17580.9 21789.7 24380.1 20538.6
classe 10   11870.5 12869.3 15784.8 18546.7 23050.9 25880.5 21700.9

class means/stddev for each band

class 1 (1997)
moyennes 7971.8 8324.33 8814.6 8664.09 9015.31 9538.61 9369.46
écart-type 664.494 701.418 852.888 930.655 1021.99 828.376 699.759
class 2 (79)
moyennes 8838.06 9346.59 10534 10553.6 14742.6 13295.6 11241.8
écart-type 928.77 1141.45 1300.27 1282.35 2055.49 1111.6 562.084
class 3 (112)
moyennes 8633.54 9013.94 10245.5 10476.8 17073.9 15341.5 12410.2
écart-type 544.033 607.832 665.275 762.651 1866.02 856.012 820.002
class 4 (160)
moyennes 9047.59 9501.68 10972.5 11595.5 17525.8 17096.4 13919.4
écart-type 1021.38 920.826 1201.83 1380.92 1698.46 1103.1 901.101
class 5 (317)
moyennes 9367.25 9960.64 11704.8 12822.7 18438.8 18797.2 15480
écart-type 482.509 593.27 836.894 1095.25 1717.77 1167.71 1123.66
class 6 (554)
moyennes 9793.56 10469.1 12494.7 14151.3 18527.5 20428.5 17087.1
écart-type 1232.55 1126.06 1387.45 1527.03 1152.56 1376.16 1276.02
class 7 (721)
moyennes 10149 10959.6 13218.6 15238.5 19468.2 21901 18407.8
écart-type 694.258 799.14 1002.51 1094.72 1057.17 1235.27 1204.21
class 8 (879)
moyennes 10545.3 11486 14062.9 16455.2 20553.2 23332.9 19668.7
écart-type 984.995 1017.89 1160.63 1227.05 904.731 1350.13 1307.15
class 9 (1009)
moyennes 11005.5 12070.7 14926 17657.2 21691.8 24558.1 20687.2
écart-type 1310.11 1314.42 1420.25 1496.34 990.198 1767.07 1709.62
class 10 (1198)
moyennes 12075.9 13287.7 16536.6 19562.6 23453.7 26286.2 22286.7
écart-type 2410.39 2280.5 2317.12 2437.84 1445.21 2800.03 2543.88

Distribution des classes
       1997         79        112        160        317
        554        721        879       1009       1198
######## iteration 1 ###########
10 classes, 91.33% points stable
Distribution des classes
       1963         97        134        167        320
        559        718        883       1246        939
######## iteration 2 ###########
10 classes, 93.25% points stable
Distribution des classes
       1889        147        157        178        325
        562        722        942       1363        741
######## iteration 3 ###########
10 classes, 93.11% points stable
Distribution des classes
       1795        232        153        188        332
        569        735       1050       1401        571
######## iteration 4 ###########
10 classes, 92.39% points stable
Distribution des classes
       1656        367        148        197        327
        581        768       1146       1420        416
######## iteration 5 ###########
10 classes, 93.13% points stable
Distribution des classes
       1561        460        145        203        323
        588        815       1217       1472        242
######## iteration 6 ###########
10 classes, 93.58% points stable
Distribution des classes
       1489        531        143        207        322
        599        856       1328       1449        102
######## iteration 7 ###########
10 classes, 94.46% points stable
Distribution des classes
       1408        611        144        207        327
        617        890       1450       1309         63
######## iteration 8 ###########
10 classes, 95.30% points stable
Distribution des classes
       1335        682        146        207        333
        636        934       1535       1160         58
######## iteration 9 ###########
10 classes, 96.23% points stable
Distribution des classes
       1276        741        146        207        339
        663        964       1576       1056         58
######## iteration 10 ###########
10 classes, 96.77% points stable
Distribution des classes
       1245        772        146        207        344
        685       1008       1598        963         58
######## iteration 11 ###########
10 classes, 97.27% points stable
Distribution des classes
       1218        799        146        206        354
        710       1031       1605        899         58
######## iteration 12 ###########
10 classes, 97.42% points stable
Distribution des classes
       1209        807        146        207        369
        730       1065       1583        852         58
######## iteration 13 ###########
10 classes, 97.65% points stable
Distribution des classes
       1207        809        144        208        390
        752       1088       1551        819         58
######## iteration 14 ###########
10 classes, 98.04% points stable
Distribution des classes
       1206        810        144        205        406
        782       1093       1531        791         58
########## final results #############
10 classes (convergence=98.0%)
class separability matrix

         1     2     3     4     5     6     7     8     9    10

  1      0
  2    1.3     0
  3    2.7   1.6     0
  4    4.2   2.9   1.2     0
  5    4.3   3.2   1.8   0.9     0
  6    6.1   4.7   2.9   1.9   0.8     0
  7    6.9   5.5   3.6   2.7   1.5   0.9     0
  8    7.9   6.4   4.5   3.6   2.3   1.8   0.9     0
  9    7.2   6.1   4.6   3.9   2.7   2.4   1.7   0.9     0
 10    3.7   3.3   3.3   3.3   3.1   3.4   3.4   3.5   3.2     0

class means/stddev for each band
class 1 (1206)
moyennes 7598.36 7874.34 8289.19 8066.25 8378.23 9026.05 8945.44
écart-type 488.014 412.695 455.617 496.048 504.221 535.432 482.295
class 2 (810)
moyennes 8592.3 9077.56 9702.47 9662.97 10012.9 10354.7 10056.4
écart-type 555.933 579.827 833.801 815.258 782.94 574.021 485.734
class 3 (144)
moyennes 8419.65 8764.67 9956.65 10029.4 17000.8 14495.5 11639
écart-type 347.949 390.554 523.602 628.656 2154.94 956.988 786.574
class 4 (205)
moyennes 8978.59 9436.55 10857.1 11482.8 17268.1 17122.9 13974.6
écart-type 310.437 365.58 520.682 684.685 1330.01 830.355 745.692
class 5 (406)
moyennes 9433.84 10073.6 11912.5 13099.6 18852.8 19035.8 15632.7
écart-type 458.667 574.321 824.748 1093.27 1706.98 1089.06 1025.15
class 6 (782)
moyennes 9827.89 10542 12599.4 14405.4 18642 21135.2 17772.6
écart-type 462.397 567.891 742.744 749.071 1014.02 807.818 912.973
class 7 (1093)
moyennes 10369.8 11262.8 13727.4 15996.5 20171.8 22965 19332.4
écart-type 560.616 665.036 776.828 750.663 912.385 932.633 1083.71
class 8 (1531)
moyennes 10930.8 12011.7 14883.7 17630.8 21760.3 24918.6 20991.9
écart-type 590.219 711.586 771.312 739.725 879.16 946.925 1242.64
class 9 (791)
moyennes 12016.2 13313.6 16680.5 19711.5 23742.9 27072.7 22986.7
écart-type 954.947 1085.14 1198.47 1172.76 1248.4 1344 1639.7
class 10 (58)
moyennes 22753.3 23009.1 26619.9 30468 25914 10463.4 9514.84
écart-type 4947.33 4201.61 3400.99 3596 4488.01 2709.38 1600.94

#################### CLASSES ####################
10 classes, 98.04% points stable
######## CLUSTER END (Thu Jul 20 20:31:04 2023) ########
\end{verbatim}
\end{scriptsize}

\subsubsection[\appendixname~\thesubsubsection]{February 2023}
\begin{scriptsize}
\begin{verbatim}
#################### CLUSTER (Thu Jul 20 20:52:42 2023) ####################
Location: Tunisia
Mapset:   PERMANENT
Group:    L8_2023f
Subgroup: res_30m
 L8_2023f_01@PERMANENT
 L8_2023f_02@PERMANENT
 L8_2023f_03@PERMANENT
 L8_2023f_04@PERMANENT
 L8_2023f_05@PERMANENT
 L8_2023f_06@PERMANENT
 L8_2023f_07@PERMANENT
Result signature file: cluster_L8_2023f
Region
  North:   4105515.00  East:    732315.00
  South:   3872685.00  West:    508155.00
  Res:          30.00  Res:         30.00
  Rows:          7761  Cols:         7472  Cells: 57990192
Mask: no
Cluster parameters
  Nombre de classes initiales: 	10
 Minimum class size:           17
 Minimum class separation:     0.000000
 Percent convergence:          98.000000
 Maximum number of iterations: 30
 Row sampling interval:        77
 Col sampling interval:        74
Sample size: 7133 points

means and standard deviations for 7 bands
moyennes 9370.75 9915.72 11415.6 12402.5 15661.8 16603.7 14586
écart-type 1388.7 1689.09 2835.91 3992.37 5540.38 6545.87 5449.16

initial means for each band

classe 1    7982.05 8226.63 8579.7 8410.17 10121.4 10057.8 9136.81
classe 2    8290.65 8601.98 9209.9 9297.37 11352.6 11512.5 10347.7
classe 3    8599.25 8977.33 9840.1 10184.6 12583.8 12967.1 11558.7
classe 4    8907.85 9352.69 10470.3 11071.8 13815 14421.7 12769.6
classe 5    9216.45 9728.04 11100.5 11958.9 15046.2 15876.4 13980.5
classe 6    9525.05 10103.4 11730.7 12846.1 16277.4 17331 15191.4
classe 7    9833.65 10478.8 12360.9 13733.3 17508.6 18785.7 16402.3
classe 8    10142.2 10854.1 12991.1 14620.5 18739.8 20240.3 17613.3
classe 9    10450.8 11229.5 13621.3 15507.7 19971 21694.9 18824.2
classe 10   10759.4 11604.8 14251.5 16394.9 21202.2 23149.6 20035.1

class means/stddev for each band

class 1 (1905)
moyennes 7812.1 7994.33 7817.89 7385.64 7324.41 7424.4 7412.99
écart-type 448.698 490.3 670.474 669.717 584.119 351.378 246.793
class 2 (87)
moyennes 8007.89 8306.43 9348.11 9482.2 12941.4 11046.7 9681.31
écart-type 1111.85 1003.77 1341.18 1952.2 1492.15 1465.13 1172.5
class 3 (183)
moyennes 8348.46 8622.51 9600.11 9642.42 15491.5 12246.9 10321.1
écart-type 855.199 933.676 1310.54 1929.9 2562.2 1382.03 1167.7
class 4 (361)
moyennes 8419.62 8706.17 9838.2 9807.6 17619.8 13894.6 11331.2
écart-type 552.509 634.89 797.882 1307.91 3294.34 1231.93 1315.07
class 5 (494)
moyennes 8823.67 9228.35 10612.9 11019.2 17393.4 15808.3 13059.3
écart-type 529.922 584.346 733.208 1191.79 2935.26 1113.77 1323.91
class 6 (656)
moyennes 9318.64 9846.84 11465.8 12397.4 17013.2 17575.4 14898.1
écart-type 551.213 624.237 781.187 1089.95 1945.1 1011.06 1126.89
class 7 (693)
moyennes 9677.52 10299.9 12164 13476.3 17561.8 19081.7 16410
écart-type 513.164 605.983 780.583 970.623 1419.38 885.72 980.699
class 8 (688)
moyennes 10018.9 10724.7 12866 14505.4 18314.6 20603.9 17854.9
écart-type 515.334 574.112 719.388 783.877 1115.33 825.965 901.936
class 9 (620)
moyennes 10379.8 11178.6 13607.6 15574.4 19242.7 21989 19138.9
écart-type 497.056 572.995 696.392 754.948 938.483 703.117 852.842
class 10 (1446)
moyennes 11195.5 12165 15166.1 17784.6 21429.4 24689.8 21681.5
écart-type 793.115 913.686 1193.57 1409.87 1424.42 1585.95 1558.53

Distribution des classes
       1905         87        183        361        494
        656        693        688        620       1446
######## iteration 1 ###########
10 classes, 84.90% points stable
Distribution des classes
       1867        157        166        359        456
        668        712        675        954       1119
######## iteration 2 ###########
10 classes, 93.06% points stable
Distribution des classes
       1855        142        255        310        461
        636        724        768       1015        967
######## iteration 3 ###########
10 classes, 95.30% points stable
Distribution des classes
       1843        121        320        286        463
        613        758        824       1033        872
######## iteration 4 ###########
10 classes, 96.50% points stable
Distribution des classes
       1844         97        346        280        470
        608        779        881       1031        797
######## iteration 5 ###########
10 classes, 96.86% points stable
Distribution des classes
       1859         69        350        276        484
        607        813        916       1008        751
######## iteration 6 ###########
10 classes, 97.39% points stable
Distribution des classes
       1861         61        346        273        495
        616        843        935        995        708
######## iteration 7 ###########
10 classes, 98.01% points stable
Distribution des classes
       1862         62        334        273        504
        627        867        943        992        669
########## final results #############
10 classes (convergence=98.0%)
class separability matrix

         1     2     3     4     5     6     7     8     9    10

  1      0
  2    1.7     0
  3    2.6   1.9     0
  4    3.1   2.1   0.8     0
  5    4.6   2.6   1.0   1.2     0
  6    5.5   3.0   1.7   2.0   0.8     0
  7    6.9   3.8   2.4   2.6   1.5   0.7     0
  8    8.3   4.5   3.2   3.4   2.4   1.5   0.8     0
  9    9.7   5.2   4.1   4.2   3.3   2.4   1.8   1.0     0
 10    8.8   5.2   4.4   4.5   3.8   3.1   2.5   1.9   1.1     0

class means/stddev for each band
class 1 (1862)
moyennes 7790.6 7968.39 7745.91 7299.51 7273.44 7407.76 7403.66
écart-type 413.971 438.626 418.601 236.773 405.936 286.15 217.916
class 2 (62)
moyennes 9994.29 10517.6 12939.6 14175.5 11509.4 7829 7594.08
écart-type 1273.62 1471.67 1722.8 2743.58 3777.99 1144.9 547.632
class 3 (334)
moyennes 8217.69 8516.43 9423.33 9501.46 14243.4 12855.8 10932.1
écart-type 629.325 589.892 597.882 777.435 1620.74 1300.66 1120.05
class 4 (273)
moyennes 8110.77 8330.43 9578.4 8995.98 21838.2 13454 10414.6
écart-type 324.184 368.487 620.97 794.986 2236.78 1423.41 1129.01
class 5 (504)
moyennes 8802.31 9204.33 10557.5 10957.7 17350.8 15909.8 13149.4
écart-type 395.53 421.184 570.838 775.511 1722.74 985.517 881.661
class 6 (627)
moyennes 9354.67 9885.15 11458.2 12498.6 15895.2 17519.6 14968.1
écart-type 509.394 570.25 760.845 923.818 1174.69 956.115 967.53
class 7 (867)
moyennes 9695.64 10329.7 12226.8 13564.7 17635.3 19215.3 16539.1
écart-type 499.684 572.239 745.313 891.507 1271.69 837.088 890.175
class 8 (943)
moyennes 10171.9 10909.8 13160.4 14943.6 18621.3 21211.3 18431.8
écart-type 484.876 555.114 712.716 798.848 1100.11 850.096 916.15
class 9 (992)
moyennes 10720 11595.5 14294.9 16621.4 20311.8 23359.6 20431.1
écart-type 499.614 567.528 687.19 781.43 950.153 890.969 995.569
class 10 (669)
moyennes 11693 12761.6 16059.1 18934.7 22466.2 25935.3 22855.3
écart-type 797.575 901.136 1056.26 1112.81 1230.5 1364.45 1309.76

#################### CLASSES ####################

10 classes, 98.01% points stable
######## CLUSTER END (Thu Jul 20 20:52:42 2023) ########
\end{verbatim}
\end{scriptsize}

\subsubsection[\appendixname~\thesubsubsection]{April 2023}
\begin{scriptsize}
\begin{verbatim}
#################### CLUSTER (Thu Jul 20 21:19:19 2023) ####################
Location: Tunisia
Mapset:   PERMANENT
Group:    L8_2023a
Subgroup: res_30m
 L8_2023a_01@PERMANENT
 L8_2023a_02@PERMANENT
 L8_2023a_03@PERMANENT
 L8_2023a_04@PERMANENT
 L8_2023a_05@PERMANENT
 L8_2023a_06@PERMANENT
 L8_2023a_07@PERMANENT
Result signature file: cluster_L8_2023a
Region
  North:   4105515.00  East:    732315.00
  South:   3872685.00  West:    508155.00
  Res:          30.00  Res:         30.00
  Rows:          7761  Cols:         7472  Cells: 57990192
Mask: no

Cluster parameters
  Nombre de classes initiales: 	10
 Minimum class size:           17
 Minimum class separation:     0.000000
 Percent convergence:          98.000000
 Maximum number of iterations: 30
 Row sampling interval:        77
 Col sampling interval:        74
Sample size: 7145 points

means and standard deviations for 7 bands
moyennes 9624.12 10209.4 11868.9 13077.9 16547 17742.1 15591.4
écart-type 1715.04 2034.42 3243.09 4425.29 6007.12 7152.5 5948.83

initial means for each band
classe 1    7909.07 8174.94 8625.83 8652.58 10539.9 10589.6 9642.56
classe 2    8290.2 8627.03 9346.51 9635.98 11874.8 12179 10964.5
classe 3    8671.32 9079.13 10067.2 10619.4 13209.7 13768.5 12286.5
classe 4    9052.44 9531.22 10787.9 11602.8 14544.6 15357.9 13608.4
classe 5    9433.56 9983.31 11508.6 12586.2 15879.5 16947.3 14930.4
classe 6    9814.68 10435.4 12229.3 13569.6 17214.4 18536.8 16252.4
classe 7    10195.8 10887.5 12950 14553 18549.4 20126.2 17574.3
classe 8    10576.9 11339.6 13670.6 15536.4 19884.3 21715.7 18896.3
classe 9    10958 11791.7 14391.3 16519.8 21219.2 23305.1 20218.3
classe 10   11339.2 12243.8 15112 17503.2 22554.1 24894.6 21540.2

class means/stddev for each band

class 1 (1925)
moyennes 7716.98 7857.65 7652.01 7321.91 7312.74 7522.79 7540.21
écart-type 408.486 447.732 606.407 475.866 250.51 185.523 175.028
class 2 (40)
moyennes 8264.02 8477.05 9552.48 9910.23 14159.3 11399.1 9700.2
écart-type 620.615 570.917 1251.93 2557.41 1535.77 1625.71 1050.78
class 3 (175)
moyennes 8275.31 8524.95 9566.72 9491.59 17158 13288.7 10884.4
écart-type 388.08 415.241 623.164 1103.41 2946.95 973.128 956.498
class 4 (340)
moyennes 8531.51 8869.09 10147.3 10248.1 18584.1 14851.9 12090.7
écart-type 440.863 528.561 682.402 1056.25 3001.74 1208.32 1333.3
class 5 (404)
moyennes 9111.06 9570.36 11135.4 11757.3 18071.7 16903.6 13979.6
écart-type 704.326 770.342 981.289 1334.51 2007.44 1192.64 1189.32
class 6 (542)
moyennes 9615.63 10188.9 11991.7 13095.7 18029.5 18776.8 15871.3
écart-type 747.281 806.356 1036.18 1298.58 1605.3 1230.38 1256.39
class 7 (758)
moyennes 9992.24 10665.9 12727.6 14268.8 18665.6 20400.5 17585.5
écart-type 1125.35 1186.15 1350.69 1490.5 1392.32 1610.71 1475.2
class 8 (824)
moyennes 10375.3 11146.7 13513.3 15400 19513.2 21992.7 19085.6
écart-type 1073.05 1119.46 1245.38 1305.87 1166.67 1469.13 1336.73
class 9 (775)
moyennes 10896.1 11771.3 14442.6 16701.1 20724.7 23377.8 20374.1
écart-type 1552.82 1560.83 1572.47 1614.67 1294.39 1970.08 1792.04
class 10 (1362)
moyennes 11578.1 12622.8 15854.1 18728.7 22688.4 26244.6 23043.9
écart-type 1007.92 1126.52 1340.64 1446.8 1377.47 1572.71 1526.63

Distribution des classes
       1925         40        175        340        404
        542        758        824        775       1362
######## iteration 1 ###########
10 classes, 89.34% points stable
Distribution des classes
       1922         60        164        308        435
        550        729        823       1056       1098
######## iteration 2 ###########
10 classes, 94.14% points stable
Distribution des classes
       1922         60        182        271        449
        545        739        901       1114        962
######## iteration 3 ###########
10 classes, 95.17% points stable
Distribution des classes
       1922         50        211        247        446
        550        756        969       1131        863
######## iteration 4 ###########
10 classes, 95.73% points stable
Distribution des classes
       1920         40        231        240        489
        477        817       1008       1122        801
######## iteration 5 ###########
10 classes, 94.75% points stable
Distribution des classes
       1915         28        252        240        589
        293        920       1051       1105        752
######## iteration 6 ###########
10 classes, 95.38% points stable
Distribution des classes
       1911         24        267        242        650
        138       1003       1106       1093        711
######## iteration 7 ###########
10 classes, 97.41% points stable
Distribution des classes
       1910         24        278        245        660
         66       1026       1169       1095        672
######## iteration 8 ###########
10 classes, 98.66% points stable
Distribution des classes
       1908         26        289        248        652
         48       1028       1199       1101        646
########## final results #############
10 classes (convergence=98.7%)
class separability matrix

         1     2     3     4     5     6     7     8     9    10

  1      0
  2    2.3     0
  3    3.9   1.9     0
  4    3.9   1.8   0.7     0
  5    5.7   2.3   1.3   1.4     0
  6    3.7   2.1   2.8   2.6   3.0     0
  7    7.2   3.0   2.3   2.3   1.0   3.4     0
  8    9.1   3.8   3.3   3.2   1.9   3.7   0.9     0
  9   11.1   4.6   4.5   4.1   2.9   4.1   1.9   1.0     0
 10   10.2   4.9   4.9   4.5   3.5   4.0   2.6   1.9   1.1     0

class means/stddev for each band

class 1 (1908)
moyennes 7701.51 7836.88 7612.22 7283.83 7301.73 7520.97 7539.59
écart-type 370.72 382.503 420.411 198.174 175.465 167.308 169.572
class 2 (26)
moyennes 9710.12 10371.3 12777.3 13564.7 10587.3 8137.38 7811.96
écart-type 692.329 873.742 1207.03 2917.97 3532.14 1547.22 857.786
class 3 (289)
moyennes 8453.18 8745.71 9779.47 9913.32 15908.7 14175 11768.9
écart-type 320.384 357.538 501.131 682.147 1465.82 1277.99 1100.35
class 4 (248)
moyennes 8232.92 8530.76 9907.34 9641.96 21574.8 14054.4 11064.1
écart-type 358.536 410.344 615.374 901.343 1920.21 1425.44 1208
class 5 (652)
moyennes 9209.56 9698.92 11275.9 12026.6 17694 17512.3 14600.1
écart-type 433.32 517.341 741.068 983.197 1679.2 971.096 978.75
class 6 (48)
moyennes 18189.2 18925.2 21660.2 24333.4 24071.2 10275.2 9130.69
écart-type 3726.57 3562.03 3270.76 3490.94 3127.04 1770.65 988.48
class 7 (1028)
moyennes 9812.47 10456.5 12442.4 13886 18351 20215 17360.5
écart-type 520.236 617.548 800.167 902.797 1349.63 1002.82 1089.65
class 8 (1199)
moyennes 10402 11206.8 13644.5 15596.9 19706.1 22436.9 19480.8
écart-type 595.273 704.018 864.903 863.605 1129.73 962.301 971.24
class 9 (1101)
moyennes 10971.6 11905.8 14817.9 17399.4 21469.6 24842.4 21762.4
écart-type 595.484 685.01 798.234 798.609 985.03 917.126 1018.93
class 10 (646)
moyennes 12099.1 13241 16729.4 19803.9 23608 27365 24063
écart-type 906.176 1029.53 1174.09 1178.84 1221.11 1263.67 1248.27

#################### CLASSES ####################
10 classes, 98.66% points stable

######## CLUSTER END (Thu Jul 20 21:19:19 2023) ########
\end{verbatim}
\end{scriptsize}

\subsubsection[\appendixname~\thesubsubsection]{July 2023}
\begin{scriptsize}
\begin{verbatim}
#################### CLUSTER (Thu Jul 20 21:30:15 2023) ####################
Location: Tunisia
Mapset:   PERMANENT
Group:    L8_2023j
Subgroup: res_30m
 L8_2023j_01@PERMANENT
 L8_2023j_02@PERMANENT
 L8_2023j_03@PERMANENT
 L8_2023j_04@PERMANENT
 L8_2023j_05@PERMANENT
 L8_2023j_06@PERMANENT
 L8_2023j_07@PERMANENT
Result signature file: cluster_L8_2023j

Region
  North:   4105515.00  East:    732315.00
  South:   3872685.00  West:    508155.00
  Res:          30.00  Res:         30.00
  Rows:          7761  Cols:         7472  Cells: 57990192
Mask: no

Cluster parameters
  Nombre de classes initiales: 	10
 Minimum class size:           17
 Minimum class separation:     0.000000
 Percent convergence:          98.000000
 Maximum number of iterations: 30
 Row sampling interval:        77
 Col sampling interval:        74
Sample size: 7142 points

means and standard deviations for 7 bands
moyennes 9968.77 10737.3 12675.1 14103.9 17263.2 19150.5 16775.2
écart-type 1645.37 1847.62 2808.35 3910.82 5323.31 6515.28 5271.38

initial means for each band
classe 1    8323.41 8889.64 9866.77 10193.1 11939.9 12635.2 11503.8
classe 2    8689.05 9300.22 10490.8 11062.2 13122.8 14083.1 12675.3
classe 3    9054.68 9710.8 11114.9 11931.3 14305.8 15530.9 13846.7
classe 4    9420.32 10121.4 11739 12800.3 15488.8 16978.7 15018.1
classe 5    9785.96 10532 12363.1 13669.4 16671.7 18426.6 16189.5
classe 6    10151.6 10942.6 12987.2 14538.5 17854.7 19874.4 17360.9
classe 7    10517.2 11353.1 13611.2 15407.6 19037.6 21322.3 18532.3
classe 8    10882.9 11763.7 14235.3 16276.6 20220.6 22770.1 19703.8
classe 9    11248.5 12174.3 14859.4 17145.7 21403.6 24217.9 20875.2
classe 10   11614.1 12584.9 15483.5 18014.8 22586.5 25665.8 22046.6

class means/stddev for each band

class 1 (1949)
moyennes 8218.34 8646.09 9100.12 8964.82 9223.69 9564.11 9359.52
écart-type 851.689 825.827 802.148 819.545 815.492 597.242 499.164
class 2 (61)
moyennes 8437.43 8885.48 10139.9 10231.5 15832.7 14110.2 11557.8
écart-type 383.839 478.586 730.507 943.428 1368.24 946.583 779.374
class 3 (137)
moyennes 8740.42 9227.01 10555.7 10895.8 16684.5 15643.8 12857.2
écart-type 424.949 408.471 539.707 1056.79 1714.19 937.839 896.209
class 4 (171)
moyennes 9073.98 9617.89 11109.6 11684.8 17390.7 17257.8 14256.6
écart-type 368.5 377.639 478.076 622.647 1678.19 681.844 835.665
class 5 (305)
moyennes 9519.02 10185.3 11958.9 12982.5 17923.7 18804.9 15761
écart-type 788.326 811.749 938.073 1059.65 1551.47 1017.86 1019.21
class 6 (607)
moyennes 9903.76 10641.4 12596.7 13979.9 18418.8 20344 17228
écart-type 656.133 688.834 851.815 964.044 1337.23 1028.37 1022.82
class 7 (891)
moyennes 10252.3 11091.8 13308.4 15062.5 19151.8 21776 18624.6
écart-type 781.367 823.812 1001.38 1094.27 1172.82 1215.35 1152.03
class 8 (938)
moyennes 10665.4 11564.5 14066.7 16176.1 20124.4 22973.6 19806.6
écart-type 1066.04 1073.63 1265.5 1403.12 1095.46 1484.2 1417.06
class 9 (819)
moyennes 11123.7 12091.3 14900 17350.7 21325.5 24138.1 20901.7
écart-type 1631.23 1528.62 1705.57 1968.71 1261.61 2087.19 1879.72
class 10 (1264)
moyennes 11670.5 12804.3 16041.2 18903.1 22973.1 26402.2 23027.3
écart-type 1202.28 1223.73 1387.16 1506.95 1238.07 1718.2 1614.06

Distribution des classes
       1949         61        137        171        305
        607        891        938        819       1264
######## iteration 1 ###########
10 classes, 91.91% points stable
Distribution des classes
       1942         74        136        182        321
        597        883        909       1065       1033
######## iteration 2 ###########
10 classes, 95.02% points stable
Distribution des classes
       1941         73        147        189        334
        594        859       1012       1056        937
######## iteration 3 ###########
10 classes, 95.25% points stable
Distribution des classes
       1940         74        158        192        335
        603        892       1109        939        900
######## iteration 4 ###########
10 classes, 94.97% points stable
Distribution des classes
       1940         74        168        191        341
        623        943       1193        763        906
######## iteration 5 ###########
10 classes, 92.92% points stable
Distribution des classes
       1940         75        173        194        351
        630       1022       1315        454        988
######## iteration 6 ###########
10 classes, 93.11% points stable
Distribution des classes
       1940         79        177        193        369
        634       1107       1408        186       1049
######## iteration 7 ###########
10 classes, 94.82% points stable
Distribution des classes
       1941         83        183        193        381
        675       1172       1397         88       1029
######## iteration 8 ###########
10 classes, 95.46% points stable
Distribution des classes
       1942         86        187        189        412
        706       1210       1371         63        976
######## iteration 9 ###########
10 classes, 96.25% points stable
Distribution des classes
       1942         94        192        175        447
        735       1233       1339         59        926
######## iteration 10 ###########
10 classes, 96.44% points stable
Distribution des classes
       1943        103        201        153        480
        766       1236       1333         58        869
######## iteration 11 ###########
10 classes, 96.43% points stable
Distribution des classes
       1943        110        214        137        503
        794       1249       1325         58        809
######## iteration 12 ###########
10 classes, 97.06% points stable
Distribution des classes
       1943        119        217        133        520
        813       1267       1309         58        763
######## iteration 13 ###########
10 classes, 97.41% points stable
Distribution des classes
       1943        130        216        132        531
        852       1255       1297         58        728
######## iteration 14 ###########
10 classes, 97.33% points stable
Distribution des classes
       1943        137        215        139        551
        876       1245       1284         58        694
######## iteration 15 ###########
10 classes, 97.44% points stable
Distribution des classes
       1943        143        219        146        554
        918       1234       1255         58        672
######## iteration 16 ###########
10 classes, 97.44% points stable
Distribution des classes
       1943        149        224        160        556
        953       1213       1235         58        651
######## iteration 17 ###########
10 classes, 98.00% points stable
Distribution des classes
       1943        154        225        170        572
        970       1193       1219         58        638
######## iteration 18 ###########
10 classes, 98.01% points stable
Distribution des classes
       1943        161        224        183        578
        991       1183       1197         58        624
########## final results #############
10 classes (convergence=98.0%)
class separability matrix

         1     2     3     4     5     6     7     8     9    10

  1      0
  2    1.9     0
  3    3.5   1.0     0
  4    3.7   1.4   0.9     0
  5    4.6   2.0   1.3   0.7     0
  6    5.9   2.8   2.2   1.2   0.8     0
  7    6.2   3.3   2.8   1.7   1.5   0.8     0
  8    7.4   4.2   3.8   2.5   2.4   1.8   0.9     0
  9    4.1   3.5   3.7   3.1   3.4   3.6   3.4   3.6     0
 10    7.3   4.5   4.3   3.0   3.1   2.6   1.7   1.0   3.4     0

class means/stddev for each band

class 1 (1943)
moyennes 8217.89 8645.61 9097.65 8962.12 9206.66 9554.53 9356.04
écart-type 852.883 826.89 801.996 819.148 753.678 572.149 495.241
class 2 (161)
moyennes 8557.61 9018.9 10320.1 10501.1 16730.8 14766.3 12023.2
écart-type 440.435 448.85 652.724 1120.75 1945.97 1168.94 926.607
class 3 (224)
moyennes 9094.89 9642.28 11100.8 11735.8 16673.1 17238.2 14383.7
écart-type 344.155 368.682 512.397 653.187 1115.15 830.197 799.113
class 4 (183)
moyennes 9287.24 10022.4 12020.2 12881.2 20658.1 18839.4 15269.1
écart-type 562.12 572.046 800.588 1118.67 1394.11 1173.62 1120.43
class 5 (578)
moyennes 9821.64 10526.5 12392 13676.4 17770.4 19849.3 16831.8
écart-type 510.795 567.389 748.082 831.169 1014.19 881.407 791.449
class 6 (991)
moyennes 10118.4 10929.5 13073.5 14804 18868.6 21833.7 18674.9
écart-type 424.861 485.242 637.679 625.63 944.455 773.539 803.384
class 7 (1183)
moyennes 10633.4 11557.6 14088.9 16201.8 20182.6 23211.5 20035.2
écart-type 581.801 673.417 788.428 738.287 915.109 932.856 955.14
class 8 (1197)
moyennes 11046.3 12075 15017.8 17615.6 21780.8 25221.6 21922
écart-type 549.329 629.431 768.178 746.48 867.908 834.515 958.731
class 9 (58)
moyennes 18960.1 19394.3 22992.3 26883.4 26181.4 12561 10794.6
écart-type 2849.9 2250.8 2134.59 2990.86 2142.15 2371.23 1248.29
class 10 (624)
moyennes 12034.3 13260.2 16675.7 19635.8 23654 27298 23836.2
écart-type 821.217 941.844 1067.11 1035.96 1080.86 1141.28 1168.92

#################### CLASSES ####################

10 classes, 98.01% points stable
######## CLUSTER END (Thu Jul 20 21:30:15 2023) ########
\end{verbatim}
\end{scriptsize}

\subsection[\appendixname~\thesubsection]{Gulf of Gabès area}

\subsubsection[\appendixname~\thesubsubsection]{February 2017}
\begin{scriptsize}
\begin{verbatim}
#################### CLUSTER (Fri Jul 21 13:52:08 2023) ####################
Location: Tunisia
Mapset:   PERMANENT
Group:    L8_2017fG
Subgroup: res_30m
 L8_2017fG_01@PERMANENT
 L8_2017fG_02@PERMANENT
 L8_2017fG_03@PERMANENT
 L8_2017fG_04@PERMANENT
 L8_2017fG_05@PERMANENT
 L8_2017fG_06@PERMANENT
 L8_2017fG_07@PERMANENT
Result signature file: cluster_L8_2017fG

Region
  North:   3946545.00  East:    699615.00
  South:   3872685.00  West:    508155.00
  Res:          30.00  Res:         30.00
  Rows:          2462  Cols:         6382  Cells: 15712484
Mask: no

Cluster parameters
  Nombre de classes initiales: 	10
 Minimum class size:           17
 Minimum class separation:     0.000000
 Percent convergence:          98.000000
 Maximum number of iterations: 30
 Row sampling interval:        24
 Col sampling interval:        63
Sample size: 7256 points

means and standard deviations for 7 bands
moyennes 9600.98 10274.8 12607 14395.7 19413.2 19976.7 17063.7
écart-type 814.942 961.157 1499.86 2267.81 2679.51 3499.3 3233.8

initial means for each band
classe 1    8786.04 9313.67 11107.1 12127.8 16733.7 16477.4 13829.9
classe 2    8967.14 9527.26 11440.4 12631.8 17329.1 17255 14548.5
classe 3    9148.24 9740.85 11773.7 13135.8 17924.5 18032.6 15267.1
classe 4    9329.34 9954.44 12107 13639.7 18520 18810.3 15985.7
classe 5    9510.43 10168 12440.3 14143.7 19115.4 19587.9 16704.4
classe 6    9691.53 10381.6 12773.6 14647.6 19710.9 20365.5 17423
classe 7    9872.63 10595.2 13106.9 15151.6 20306.3 21143.1 18141.6
classe 8    10053.7 10808.8 13440.2 15655.5 20901.8 21920.7 18860.2
classe 9    10234.8 11022.4 13773.5 16159.5 21497.2 22698.4 19578.8
classe 10   10415.9 11236 14106.8 16663.5 22092.7 23476 20297.5

class means/stddev for each band

class 1 (896)
moyennes 8420.36 8876.24 10308.6 10692.7 15650.2 13327.2 11163.1
écart-type 732.041 847.201 1353.01 1949.41 4720.36 3398.33 2382.42
class 2 (428)
moyennes 8883.83 9445.33 11273.3 12267.8 18305.2 17187 14226.4
écart-type 334.897 370.542 569.166 823.934 1672.12 579.725 900.696
class 3 (578)
moyennes 9107.33 9682.45 11631.9 12810.1 18585.6 18054.8 15047.7
écart-type 322.719 366.317 536.513 691.263 1462.37 521.216 808.101
class 4 (680)
moyennes 9252.08 9878.84 11969.9 13395.6 18891.1 18897.8 15895.3
écart-type 269.631 316.215 466.562 560.909 1333.12 501.987 792.521
class 5 (774)
moyennes 9471.44 10124 12360.1 14007.8 19173.9 19710.1 16686.4
écart-type 292.526 354.697 510.348 563.817 1173.13 479.612 778.929
class 6 (787)
moyennes 9645.81 10325.2 12692.5 14551.9 19566.6 20533.6 17489.8
écart-type 294.813 349.821 468.087 485.626 1107.43 471.231 755.688
class 7 (800)
moyennes 9831.55 10546.5 13049.5 15150.7 19979.8 21320.5 18296.2
écart-type 299.623 376.498 490.07 470.768 986.948 489.449 763.155
class 8 (730)
moyennes 10034.6 10786.6 13447.8 15759.2 20422.4 22059.3 19025.7
écart-type 367.461 446.496 522.651 461.056 909.141 437.263 718.754
class 9 (561)
moyennes 10191.8 10960.6 13757.8 16334 20839 22811.4 19837.5
écart-type 403.7 500.265 598.662 485.404 812.378 464.517 779.554
class 10 (1022)
moyennes 10696.8 11567.6 14698.3 17639.7 22107.6 24457 21411.2
écart-type 682.345 807.334 992.89 1022.63 1135.64 1178.93 1293.71

Distribution des classes
        896        428        578        680        774
        787        800        730        561       1022
######## iteration 1 ###########
10 classes, 81.48% points stable
Distribution des classes
        491        851        562        683        787
        777        781        726        840        758
######## iteration 2 ###########
10 classes, 89.71% points stable
Distribution des classes
        301        868        739        680        791
        785        778        783        921        610
######## iteration 3 ###########
10 classes, 87.72% points stable
Distribution des classes
        242        733        886        724        776
        822        783        842        922        526
######## iteration 4 ###########
10 classes, 85.56% points stable
Distribution des classes
        225        612        922        769        830
        804        873        817        925        479
######## iteration 5 ###########
10 classes, 92.56% points stable
Distribution des classes
        217        544        892        791        905
        816        916        791        927        457
######## iteration 6 ###########
10 classes, 96.20% points stable
Distribution des classes
        215        487        861        804        954
        828        963        775        923        446
######## iteration 7 ###########
10 classes, 96.98% points stable
Distribution des classes
        214        440        833        804        987
        854        996        773        916        439
######## iteration 8 ###########
10 classes, 97.53% points stable
Distribution des classes
        212        404        806        820       1002
        862       1024        776        916        434
######## iteration 9 ###########
10 classes, 97.93% points stable
Distribution des classes
        212        376        777        820       1031
        871       1039        783        917        430
######## iteration 10 ###########
10 classes, 98.14% points stable
Distribution des classes
        212        358        750        815       1053
        873       1058        791        916        430
########## final results #############
10 classes (convergence=98.1%)
class separability matrix

         1     2     3     4     5     6     7     8     9    10

  1      0
  2    1.6     0
  3    2.7   0.9     0
  4    3.6   1.2   0.7     0
  5    4.2   1.8   1.0   0.9     0
  6    4.6   2.1   1.5   1.3   0.8     0
  7    5.3   2.6   2.0   1.9   1.2   0.8     0
  8    5.4   2.8   2.4   2.3   1.8   1.1   0.9     0
  9    6.0   3.2   2.8   2.9   2.4   1.8   1.3   0.8     0
 10    5.5   3.2   3.1   3.0   2.8   2.3   2.0   1.4   1.0     0

class means/stddev for each band

class 1 (212)
moyennes 7813.15 8197.51 9084.63 8816.3 8243.97 7935.18 7755.36
écart-type 814.555 947.359 1574.72 1862.92 1596.3 1171.78 809.17
class 2 (358)
moyennes 8356.64 8782.41 10317.7 10511.9 19280.2 14010.4 11241.8
écart-type 606.641 701.032 1023.74 1524.8 2672.14 1678.93 1298.85
class 3 (750)
moyennes 8997.62 9552.15 11319.9 12457.4 17011.6 16888.3 14066.3
écart-type 415.049 486.671 798.943 1077.25 1125.04 1014.23 1219.72
class 4 (815)
moyennes 9016.65 9642.59 11729.4 12736.2 20014 18064 14804
écart-type 311.118 356.744 520.475 749.768 964.662 768.236 802.307
class 5 (1053)
moyennes 9392.59 9983.83 12042.2 13742.1 18114.9 19491.5 16715.9
écart-type 270.12 337.547 521.589 605.423 787.506 654.373 746.524
class 6 (873)
moyennes 9672.07 10413.1 12873.8 14503.7 20397 20217.2 16884.9
écart-type 333.809 375.527 502.327 621.055 937.099 592.942 706.649
class 7 (1058)
moyennes 9741.79 10392 12812.5 15136.3 19322.7 21561.2 18828.7
écart-type 268.955 337.404 435.648 518.736 644.368 596.961 682.69
class 8 (791)
moyennes 10260.6 11093 13876.6 16017.5 21053.1 22018.9 18661.2
écart-type 397.939 468.026 563.151 628.521 947.386 587.276 686.455
class 9 (916)
moyennes 10250.3 11017.7 13897.9 16737 21058.3 23511.8 20696.2
écart-type 416.333 492.59 591.411 567.249 784.242 592.674 763.229
class 10 (430)
moyennes 11101.4 12056.4 15407.6 18493.3 22963.9 25406.4 22205.8
écart-type 766.872 901.189 1028.45 949.159 1010.83 1150.27 1314.62

#################### CLASSES ####################

10 classes, 98.14% points stable

######## CLUSTER END (Fri Jul 21 13:52:08 2023) ########

\end{verbatim}
\end{scriptsize}

\subsubsection[\appendixname~\thesubsubsection]{April 2017}
\begin{scriptsize}
\begin{verbatim}
#################### CLUSTER (Fri Jul 21 14:10:15 2023) ####################
Location: Tunisia_Gabes
Mapset:   PERMANENT
Group:    L8_2017a
Subgroup: res_30m
 L8_2017a_01@PERMANENT
 L8_2017a_02@PERMANENT
 L8_2017a_03@PERMANENT
 L8_2017a_04@PERMANENT
 L8_2017a_05@PERMANENT
 L8_2017a_06@PERMANENT
 L8_2017a_07@PERMANENT
Result signature file: cluster_L8_2017a

Region
  North:   3946515.00  East:    699315.00
  South:   3713085.00  West:    470385.00
  Res:          30.00  Res:         30.00
  Rows:          7781  Cols:         7631  Cells: 59376811
Mask: no

Cluster parameters
  Nombre de classes initiales: 	10
 Minimum class size:           17
 Minimum class separation:     0.000000
 Percent convergence:          98.000000
 Maximum number of iterations: 30
 Row sampling interval:        77
 Col sampling interval:        76
Sample size: 6930 points

means and standard deviations for 7 bands
moyennes 10096.8 11006.2 13497.3 15453.3 19400.5 20271.5 17762.8
écart-type 1651.47 1883 3004.55 4404.22 6030.12 6481.46 5398.28

initial means for each band
classe 1    8445.31 9123.19 10492.7 11049.1 13370.4 13790.1 12364.5
classe 2    8812.31 9541.64 11160.4 12027.8 14710.4 15230.4 13564.1
classe 3    9179.3 9960.08 11828.1 13006.5 16050.5 16670.7 14763.8
classe 4    9546.29 10378.5 12495.7 13985.2 17390.5 18111.1 15963.4
classe 5    9913.29 10797 13163.4 14964 18730.5 19551.4 17163
classe 6    10280.3 11215.4 13831.1 15942.7 20070.6 20991.7 18362.6
classe 7    10647.3 11633.9 14498.8 16921.4 21410.6 22432 19562.2
classe 8    11014.3 12052.3 15166.5 17900.1 22750.6 23872.4 20761.8
classe 9    11381.3 12470.8 15834.1 18878.8 24090.6 25312.7 21961.5
classe 10   11748.3 12889.2 16501.8 19857.5 25430.7 26753 23161.1

class means/stddev for each band

class 1 (1324)
moyennes 7449.67 8045.93 8391.75 7682.48 7717.8 7967.35 7932.56
écart-type 601.978 688.822 1056.25 1103.94 1514.98 743.361 454.946
class 2 (45)
moyennes 8567.09 9088.78 10644.3 10907.2 19039.7 14143.6 11551.5
écart-type 734.818 906.559 1206.74 1840.83 3137.49 1455.69 1367.33
class 3 (103)
moyennes 9406.83 10051.4 11965.6 12884.7 18767.6 16076.7 13250.3
écart-type 1534.45 1613.8 1806.2 2304.16 2069.48 1990.48 1761.81
class 4 (231)
moyennes 9362.89 10057.8 12124.5 13338.4 19262 18109.5 15161.6
écart-type 883.865 959.948 1134.32 1445.66 1634.3 1303.01 1417.79
class 5 (440)
moyennes 9684.4 10454.2 12748.4 14396.1 19558.7 19885.7 16945.7
écart-type 952.527 1026.91 1146.44 1299.52 1166.05 1174.09 1355.64
class 6 (824)
moyennes 10005.2 10867.6 13447 15484.6 20444.2 21408.4 18408.2
écart-type 565.871 730.278 963.802 1086.09 1046.68 913.827 1239.35
class 7 (1073)
moyennes 10391.9 11352.3 14229 16647.9 21522.7 22742.9 19704.7
écart-type 535.362 701.613 913.6 954.939 915.697 892.842 1162.53
class 8 (1186)
moyennes 10830.6 11865.6 15021 17843.3 22606 24037.1 20923.5
écart-type 582.107 739.896 933.071 935.928 865.846 876.463 1252.46
class 9 (946)
moyennes 11401.6 12519.6 15969.5 19150.6 23774.4 25206.2 21965.3
écart-type 618.501 780.555 962.566 933.771 863.111 882.971 1380.88
class 10 (758)
moyennes 12273.3 13457.2 17194.9 20825.1 25252 26795.4 23541.6
écart-type 941.477 1046 1126.32 1072.55 966.14 1056.43 1504.96

Distribution des classes
       1324         45        103        231        440
        824       1073       1186        946        758
######## iteration 1 ###########
10 classes, 91.76% points stable
Distribution des classes
       1287         88        106        238        453
        820       1044       1169       1048        677
######## iteration 2 ###########
10 classes, 95.44% points stable
Distribution des classes
       1281         92         95        287        450
        802       1043       1162       1090        628
######## iteration 3 ###########
10 classes, 95.31% points stable
Distribution des classes
       1279         91        118        288        472
        776       1048       1165       1089        604
######## iteration 4 ###########
10 classes, 95.17% points stable
Distribution des classes
       1280         93        141        281        498
        787       1032       1152       1071        595
######## iteration 5 ###########
10 classes, 94.88% points stable
Distribution des classes
       1281         94        160        280        519
        820       1009       1130       1038        599
######## iteration 6 ###########
10 classes, 95.17% points stable
Distribution des classes
       1280         94        173        283        553
        866        972       1117        973        619
######## iteration 7 ###########
10 classes, 96.05% points stable
Distribution des classes
       1278         95        171        293        593
        907        952       1085        911        645
######## iteration 8 ###########
10 classes, 96.81% points stable
Distribution des classes
       1276         96        176        298        622
        947        940       1064        852        659
######## iteration 9 ###########
10 classes, 97.55% points stable
Distribution des classes
       1276         95        180        302        661
        960        929       1050        821        656
######## iteration 10 ###########
10 classes, 98.08% points stable
Distribution des classes
       1275         97        183        308        677
        982        932       1025        800        651
########## final results #############
10 classes (convergence=98.1%)

class separability matrix

         1     2     3     4     5     6     7     8     9    10

  1      0
  2    1.6     0
  3    4.9   1.5     0
  4    4.7   1.0   1.7     0
  5    7.9   1.8   1.5   1.0     0
  6   10.1   2.3   1.5   1.7   0.9     0
  7   10.4   2.5   1.1   2.2   1.6   1.0     0
  8   11.6   3.0   1.4   2.7   2.2   1.4   0.8     0
  9   10.8   2.9   1.1   3.0   2.6   2.2   1.2   1.1     0
 10   11.1   3.4   1.5   3.5   3.1   2.7   1.9   1.5   0.8     0

class means/stddev for each band

class 1 (1275)
moyennes 7384.34 7981.24 8282.87 7536.8 7457.55 7850.38 7869.24
écart-type 395.346 506.206 802.302 683.91 454.361 223.512 196.105
class 2 (97)
moyennes 9099.24 9689.85 11336.3 11755.4 16659.1 12387 10366.2
écart-type 1466.1 1630.46 2089.64 2626.92 3722.46 2206.88 1477.76
class 3 (183)
moyennes 12036.4 13374.8 16692.5 19380.7 21167.8 19219.5 15013.3
écart-type 1207.44 1208.99 1351.85 1693.57 1857.55 2764.88 2062.28
class 4 (308)
moyennes 9111.89 9772 11755.7 12758.5 19337.6 17726.5 14719
écart-type 520.922 600.98 817.461 1232.86 1902.76 1127.47 1257.82
class 5 (677)
moyennes 9627.97 10410.8 12754.5 14415.7 19892.1 20399.6 17471.7
écart-type 467.912 533.105 699.42 760.713 1312.43 769.372 830.018
class 6 (982)
moyennes 9975.51 10776.8 13391.1 15734 20625.3 22536.6 19689.2
écart-type 302.219 373.401 503.406 606.224 919.087 788.271 953.607
class 7 (932)
moyennes 10698.4 11791.4 14850.1 17206 22390.6 22895.3 19715.6
écart-type 412.746 488.688 621.075 737.899 871.665 909.288 833.656
class 8 (1025)
moyennes 10666.5 11576.2 14724.2 17934.1 22579.8 25189.7 22321
écart-type 403.426 475.312 579.777 734.066 852.102 809.157 1024.56
class 9 (800)
moyennes 11764.7 13015.9 16557.6 19525.7 24085 24550.5 21135.4
écart-type 524.151 591.308 765.54 937.183 907.841 934.036 1110.63
class 10 (651)
moyennes 12341 13518.8 17260.6 20916 25330.8 26965.2 23734.4
écart-type 927.703 1020.49 1077.05 1023.42 943.426 966.84 1316.47

#################### CLASSES ####################
10 classes, 98.08% points stable
######## CLUSTER END (Fri Jul 21 14:10:15 2023) ########
\end{verbatim}
\end{scriptsize}

\subsubsection[\appendixname~\thesubsubsection]{July 2017}
\begin{scriptsize}
\begin{verbatim}
#################### CLUSTER (Fri Jul 21 14:27:47 2023) ####################
Location: Tunisia_Gabes
Mapset:   PERMANENT
Group:    L8_2017j
Subgroup: res_30m
 L8_2017j_01@PERMANENT
 L8_2017j_02@PERMANENT
 L8_2017j_03@PERMANENT
 L8_2017j_04@PERMANENT
 L8_2017j_05@PERMANENT
 L8_2017j_06@PERMANENT
 L8_2017j_07@PERMANENT
Result signature file: cluster_L8_2017j

Region
  North:   3946515.00  East:    699315.00
  South:   3713085.00  West:    472455.00
  Res:          30.00  Res:         30.00
  Rows:          7781  Cols:         7562  Cells: 58839922
Mask: no

Cluster parameters
  Nombre de classes initiales: 	10
 Minimum class size:           17
 Minimum class separation:     0.000000
 Percent convergence:          98.000000
 Maximum number of iterations: 30
 Row sampling interval:        77
 Col sampling interval:        75
Sample size: 7022 points

means and standard deviations for 7 bands

moyennes 10695.9 11773.9 14451.8 16872.1 19854.6 21937.3 18723.9
écart-type 2208.72 2287.02 2859.57 4038.55 5236.48 6384.22 5045.35

initial means for each band

classe 1    8487.16 9486.83 11592.3 12833.5 14618.1 15553.1 13678.5
classe 2    8977.99 9995.06 12227.7 13731 15781.8 16971.8 14799.7
classe 3    9468.81 10503.3 12863.2 14628.4 16945.4 18390.5 15920.9
classe 4    9959.64 11011.5 13498.6 15525.9 18109.1 19809.3 17042.1
classe 5    10450.5 11519.7 14134.1 16423.3 19272.8 21228 18163.3
classe 6    10941.3 12028 14769.6 17320.8 20436.4 22646.7 19284.4
classe 7    11432.1 12536.2 15405 18218.2 21600.1 24065.4 20405.6
classe 8    11923 13044.4 16040.5 19115.7 22763.7 25484.1 21526.8
classe 9    12413.8 13552.6 16675.9 20013.1 23927.4 26902.8 22648
classe 10   12904.6 14060.9 17311.4 20910.6 25091.1 28321.6 23769.2

class means/stddev for each band

class 1 (1413)
moyennes 8300.83 9092.33 10062.2 9898.96 10083.5 9903.19 9437.23
écart-type 822.138 773.955 735.754 860.389 1230.39 906.692 635.83
class 2 (29)
moyennes 8781 10016.4 12242.2 13681.9 17253.4 16817.8 14172.7
écart-type 1175.78 731.914 588.73 641.197 1223.17 819.77 796.09
class 3 (49)
moyennes 9339 10419.7 12876 14256.6 18945.4 18379.7 15046.8
écart-type 797.33 617.482 557.021 818.221 1839.55 810.887 964.51
class 4 (157)
moyennes 9792.68 10765.5 13120.4 14999.7 19114.5 20140.6 16799.2
écart-type 645.038 631.69 722.368 722.24 1205.58 902.976 1054.63
class 5 (372)
moyennes 10108.2 11177.1 13789.9 16018.6 19902.3 21696.7 18177.4
écart-type 694.238 687.69 736.815 767.395 805.608 926.482 930.716
class 6 (837)
moyennes 10437.9 11524.4 14382 17035.3 20762.6 23164.8 19544.3
écart-type 686.415 751.505 918.371 1024.11 798.453 1180.95 1237.55
class 7 (1303)
moyennes 10790.9 11911.8 14948.2 17945.7 21670.8 24593.3 20783.6
écart-type 559.2 649.903 740.592 720.592 563.792 772.621 1061.84
class 8 (1382)
moyennes 11293.3 12485.3 15738.3 19047.7 22743.2 25821 21762.7
écart-type 699.675 763.774 837.092 808.479 601.699 812.81 1200.12
class 9 (950)
moyennes 11947.6 13173.5 16650.1 20263.5 23880.3 26922.5 22650.8
écart-type 793.404 837.78 867.791 804.442 621.054 803.746 1287.97
class 10 (530)
moyennes 14363.9 15552.8 18875.5 22383.4 25669.7 27800.3 23702.4
écart-type 4989.89 4817.02 3823.11 3013.73 2313.99 1611.2 1250.93

Distribution des classes
       1413         29         49        157        372
        837       1303       1382        950        530
######## iteration 1 ###########
10 classes, 90.02% points stable
Distribution des classes
       1378         67         67        163        385
        835       1266       1344       1239        278
######## iteration 2 ###########
10 classes, 93.75% points stable
Distribution des classes
       1374         69         79        172        408
        837       1217       1442       1315        109
######## iteration 3 ###########
10 classes, 95.57% points stable
Distribution des classes
       1373         74         74        193        419
        858       1202       1526       1229         74
######## iteration 4 ###########
10 classes, 95.98% points stable
Distribution des classes
       1373         75         79        218        435
        865       1222       1529       1162         64
######## iteration 5 ###########
10 classes, 95.93% points stable
Distribution des classes
       1373         74        133        222        457
        878       1245       1504       1075         61
######## iteration 6 ###########
10 classes, 96.65% points stable
Distribution des classes
       1373         72        167        226        472
        904       1290       1441       1017         60
######## iteration 7 ###########
10 classes, 97.28% points stable
Distribution des classes
       1373         70        185        225        487
        945       1304       1394        980         59
######## iteration 8 ###########
10 classes, 97.82% points stable
Distribution des classes
       1373         70        191        226        502
        955       1341       1366        942         56
######## iteration 9 ###########
10 classes, 98.05% points stable
Distribution des classes
       1373         70        199        227        517
        968       1357       1353        905         53
########## final results #############
10 classes (convergence=98.0%)
class separability matrix

         1     2     3     4     5     6     7     8     9    10

  1      0
  2    2.0     0
  3    4.2   2.3     0
  4    4.0   1.1   1.7     0
  5    6.1   2.0   1.5   0.8     0
  6    7.8   2.6   1.4   1.5   0.8     0
  7    9.1   3.2   1.3   2.2   1.6   0.9     0
  8    9.8   3.7   1.1   2.8   2.4   1.8   1.0     0
  9    9.2   4.1   1.0   3.2   3.0   2.6   2.0   1.1     0
 10    4.5   3.4   1.8   3.0   2.9   2.8   2.7   2.5   2.1     0

class means/stddev for each band

class 1 (1373)
moyennes 8282.61 9069.2 10008.5 9818.7 9926.83 9786.84 9367.77
écart-type 794.142 748.651 656.803 716.264 811.225 553.241 452.249
class 2 (70)
moyennes 8868.64 9941.03 12053.7 13083.4 16323.9 15135.8 12797.5
écart-type 1274.81 979.252 802.061 918.772 1690.31 1893.01 1558.76
class 3 (199)
moyennes 13333.6 14748.1 18276 21489.8 23954.9 22786.1 17551.7
écart-type 2020.11 1922.36 1691.73 1790.07 1760.72 2618.23 2177.75
class 4 (227)
moyennes 9770.12 10785.3 13200.3 14992.4 19217.1 19744.5 16376.9
écart-type 747.489 706.908 807.57 980.944 1395.54 1069.05 1168.08
class 5 (517)
moyennes 10204.9 11287.4 13956.8 16246 20080.2 22028.1 18455.5
écart-type 594.121 608.87 735.658 784.87 831.763 711.533 742.099
class 6 (968)
moyennes 10397.5 11466.6 14341.5 17102.4 20910.5 23744.2 20131.4
écart-type 439.108 497.793 565.592 524.119 600.661 597.504 746.382
class 7 (1357)
moyennes 10837.2 11959.7 15037.1 18139 21927.2 25098.3 21300.4
écart-type 427.147 485.901 517.103 463.152 472.522 575.635 784.406
class 8 (1353)
moyennes 11391.8 12591.4 15915.4 19351.2 23084.8 26401.5 22354.3
écart-type 470.989 508.445 523.59 502.188 483.015 601.904 855.574
class 9 (905)
moyennes 12483.8 13728.2 17323.2 21038.5 24565.1 27658.2 23311.3
écart-type 806.179 812.623 809.202 724.708 650.223 826.14 980.15
class 10 (53)
moyennes 28015 28684.9 29197.2 30338.3 31589 25426.9 23308.2
écart-type 5083.59 5012.35 4125.93 3876.02 3383.29 3377.05 2727.85

#################### CLASSES ####################
10 classes, 98.05% points stable

######## CLUSTER END (Fri Jul 21 14:27:47 2023) ########
\end{verbatim}
\end{scriptsize}

\subsubsection[\appendixname~\thesubsubsection]{February 2023}
\begin{scriptsize}
\begin{verbatim}
#################### CLUSTER (Fri Jul 21 15:26:19 2023) ####################
Location: Tunisia_Gabes
Mapset:   PERMANENT
Group:    L8_2023f
Subgroup: res_30m
 L8_2023f_01@PERMANENT
 L8_2023f_02@PERMANENT
 L8_2023f_03@PERMANENT
 L8_2023f_04@PERMANENT
 L8_2023f_05@PERMANENT
 L8_2023f_06@PERMANENT
 L8_2023f_07@PERMANENT
Result signature file: cluster_L8_2023f

Region
  North:   3946245.00  East:    696915.00
  South:   3713085.00  West:    472455.00
  Res:          30.00  Res:         30.00
  Rows:          7772  Cols:         7482  Cells: 58150104
Mask: no

Cluster parameters
  Nombre de classes initiales: 	10
 Minimum class size:           17
 Minimum class separation:     0.000000
 Percent convergence:          98.000000
 Maximum number of iterations: 30
 Row sampling interval:        77
 Col sampling interval:        74
Sample size: 7137 points

means and standard deviations for 7 bands

moyennes 10626.6 11452.5 13947.8 16471.4 19195.8 21475.5 19086.3
écart-type 1529.21 1818.12 3149.03 4816.68 6006.71 7307.63 6346.7

initial means for each band

classe 1    9097.39 9634.36 10798.7 11654.7 13189.1 14167.8 12739.6
classe 2    9437.22 10038.4 11498.5 12725.1 14523.9 15791.7 14150
classe 3    9777.04 10442.4 12198.3 13795.5 15858.7 17415.7 15560.4
classe 4    10116.9 10846.4 12898.1 14865.9 17193.6 19039.6 16970.7
classe 5    10456.7 11250.5 13597.9 15936.2 18528.4 20663.5 18381.1
classe 6    10796.5 11654.5 14297.7 17006.6 19863.2 22287.4 19791.5
classe 7    11136.3 12058.5 14997.5 18077 21198 23911.3 21201.9
classe 8    11476.2 12462.6 15697.2 19147.4 22532.9 25535.3 22612.2
classe 9    11816 12866.6 16397 20217.7 23867.7 27159.2 24022.6
classe 10   12155.8 13270.6 17096.8 21288.1 25202.5 28783.1 25433

class means/stddev for each band

class 1 (1352)
moyennes 8213.76 8527.83 8428.42 7773.63 7757.05 7633.57 7529.31
écart-type 493.475 618.147 975.023 1122.3 1965.89 1303.68 896.44
class 2 (81)
moyennes 9005.85 9589.73 11191.1 12305 16472 15910.8 13361.2
écart-type 520.676 601.107 861.095 1296.31 2148.91 878.122 1178.57
class 3 (183)
moyennes 9508.17 10135.7 12022.8 13571.4 17342.2 17589.1 14636.4
écart-type 709.102 753.728 962.064 1390.9 1695.97 851.278 1281.74
class 4 (292)
moyennes 9953.89 10683.5 12830.2 14804.4 18129.3 19188.2 16176.9
écart-type 658.523 755.555 985.022 1304.05 1121.85 885.681 1608.39
class 5 (478)
moyennes 10172.9 10950.1 13318.4 15625.1 19046.6 21044.1 18195
écart-type 670.592 743.411 936.017 1105.03 991.931 973.072 1487.04
class 6 (685)
moyennes 10484.1 11320.7 13965.7 16644.4 20135.8 22676.7 19884.5
écart-type 582.658 687.532 846.009 933.87 873.192 776.581 1285.9
class 7 (938)
moyennes 10912.7 11820 14783.9 17877.3 21266.8 24188.6 21275.4
écart-type 593.95 681.313 800.853 897.693 785.301 711.526 1339.18
class 8 (1148)
moyennes 11332.3 12309.5 15576 19081.9 22448.4 25682.4 22757
écart-type 627.297 709.767 831.739 899.35 806.057 660.271 1455.74
class 9 (1057)
moyennes 11872.1 12924 16475.4 20359.3 23611.4 27037.1 24116.5
écart-type 652.914 709.948 741.559 700.521 602.922 609.341 1055.52
class 10 (923)
moyennes 12483.7 13637.6 17552.8 21862 25067.9 28706.5 25638.5
écart-type 668.393 742.409 809.538 820.03 812.635 1055.36 1160.25

Distribution des classes
       1352         81        183        292        478
        685        938       1148       1057        923
######## iteration 1 ###########
10 classes, 94.30% points stable
Distribution des classes
       1282        163        184        291        475
        697        929       1130       1125        861
######## iteration 2 ###########
10 classes, 96.13% points stable
Distribution des classes
       1272        144        231        286        488
        688        934       1132       1130        832
######## iteration 3 ###########
10 classes, 95.80% points stable
Distribution des classes
       1270        119        278        236        531
        703        918       1144       1124        814
######## iteration 4 ###########
10 classes, 97.66% points stable
Distribution des classes
       1267        114        303        211        539
        732        899       1143       1132        797
######## iteration 5 ###########
10 classes, 98.43% points stable
Distribution des classes
       1267        110        319        193        548
        744        908       1131       1136        781
########## final results #############
10 classes (convergence=98.4%)
class separability matrix

         1     2     3     4     5     6     7     8     9    10

  1      0
  2    2.0     0
  3    5.0   1.0     0
  4    4.2   1.6   1.2     0
  5    8.6   2.0   1.1   1.1     0
  6   11.5   2.7   2.0   1.2   1.0     0
  7   12.0   3.1   2.6   1.3   1.7   0.9     0
  8   15.2   3.7   3.4   1.8   2.6   1.8   0.8     0
  9   14.0   4.0   3.8   2.0   3.2   2.5   1.5   1.0     0
 10   13.5   4.4   4.3   2.3   3.8   3.2   2.3   1.8   0.8     0

class means/stddev for each band

class 1 (1267)
moyennes 8181.92 8483.49 8291 7554.19 7292.13 7326.34 7326.94
écart-type 449.862 567.051 790.694 667.195 406.902 213.103 154.831
class 2 (110)
moyennes 8729.98 9246.55 10601.2 11231 15215.6 12865.7 10973.6
écart-type 746.132 874.122 1146.97 1478.78 2959.09 2150.49 1738.29
class 3 (319)
moyennes 9380.96 10009.1 11837.6 13322.1 17422 17807.7 15092.5
écart-type 533.121 585.83 765.794 1116.66 1815.59 1124.97 1275.23
class 4 (193)
moyennes 11455.5 12433.2 15257.2 18054 19921.1 19271.9 14264
écart-type 1039.64 1248.41 1745.86 2319.62 2582.28 1810.68 1776.86
class 5 (548)
moyennes 9939.5 10685.5 12956.9 15135.1 18730.9 20971.8 18352.7
écart-type 446.917 501.422 648.422 734.099 1013.86 858.838 935.615
class 6 (744)
moyennes 10404.3 11222.8 13865.1 16559.5 20115.6 22910 20243.3
écart-type 412.41 478.223 571.966 580.807 738.136 743.743 845.179
class 7 (908)
moyennes 11002.2 11921.3 14906.2 18019.1 21377.5 24276.9 21354.1
écart-type 558.993 624.873 702.638 757.8 708.047 707.491 1076.51
class 8 (1131)
moyennes 11202.6 12161.6 15429.6 18991.1 22425.4 25891.6 23203.5
écart-type 494.125 544.211 576.822 565.421 600.769 622.103 827.791
class 9 (1136)
moyennes 12005.8 13077.9 16667.7 20568 23763.6 27103.6 24052.8
écart-type 584.933 643.183 687.9 710.662 663.637 783.741 1471.14
class 10 (781)
moyennes 12580.8 13747.7 17691.9 22018 25205.2 28832.8 25718.6
écart-type 671.681 737.767 787.03 783.652 798.234 1091.89 1209.47

#################### CLASSES ####################

10 classes, 98.43% points stable

######## CLUSTER END (Fri Jul 21 15:26:19 2023) ########
\end{verbatim}
\end{scriptsize}

\subsubsection[\appendixname~\thesubsubsection]{April 2023}
\begin{scriptsize}
\begin{verbatim}
#################### CLUSTER (Fri Jul 21 15:37:15 2023) ####################
Location: Tunisia_Gabes
Mapset:   PERMANENT
Group:    L8_2023a
Subgroup: res_30m
 L8_2023a_01@PERMANENT
 L8_2023a_02@PERMANENT
 L8_2023a_03@PERMANENT
 L8_2023a_04@PERMANENT
 L8_2023a_05@PERMANENT
 L8_2023a_06@PERMANENT
 L8_2023a_07@PERMANENT
Result signature file: cluster_L8_2023a

Region
  North:   3946245.00  East:    696915.00
  South:   3713085.00  West:    472455.00
  Res:          30.00  Res:         30.00
  Rows:          7772  Cols:         7482  Cells: 58150104
Mask: no

Cluster parameters
  Nombre de classes initiales: 	10
 Minimum class size:           17
 Minimum class separation:     0.000000
 Percent convergence:          98.000000
 Maximum number of iterations: 30
 Row sampling interval:        77
 Col sampling interval:        74

Sample size: 7140 points

means and standard deviations for 7 bands

moyennes 10857 11754.4 14441.1 17162.5 20232.2 22629.6 20028.1
écart-type 1719.16 2004.74 3394.11 5149.55 6446.38 7767.49 6667.49

initial means for each band

classe 1    9137.87 9749.62 11047 12013 13785.8 14862.1 13360.6
classe 2    9519.9 10195.1 11801.2 13157.3 15218.4 16588.2 14842.3
classe 3    9901.94 10640.6 12555.5 14301.7 16650.9 18314.3 16323.9
classe 4    10284 11086.1 13309.7 15446 18083.4 20040.4 17805.6
classe 5    10666 11531.6 14064 16590.4 19516 21766.5 19287.3
classe 6    11048 11977.1 14818.2 17734.7 20948.5 23492.6 20768.9
classe 7    11430.1 12422.6 15572.5 18879.1 22381 25218.7 22250.6
classe 8    11812.1 12868.1 16326.7 20023.4 23813.5 26944.8 23732.3
classe 9    12194.2 13313.6 17080.9 21167.7 25246.1 28670.9 25213.9
classe 10   12576.2 13759.1 17835.2 22312.1 26678.6 30397.1 26695.6

class means/stddev for each band

class 1 (1356)
moyennes 8009.57 8375.09 8293.76 7623.87 7667.25 7629.09 7563.96
écart-type 482.47 543.416 911.476 974.43 2102.8 1071.2 672.183
class 2 (52)
moyennes 9353.27 9998.56 11892.3 13037.5 18513.8 16113.4 13096.9
écart-type 1106.89 1129.38 1631.67 2474.31 2462.57 1602.12 1551.32
class 3 (115)
moyennes 9907.53 10655.8 12826.1 14407.6 19093 17760.7 14710.3
écart-type 1441.44 1470.69 1816.45 2416.66 2215.99 2011.03 1949.65
class 4 (214)
moyennes 10106.4 10911 13191.7 15117.4 19485.1 20129.2 17066.3
écart-type 1048.72 1111.75 1322.41 1592.75 1492.55 1369.62 1600.86
class 5 (429)
moyennes 10337.2 11181.2 13735.3 16091.2 20286.9 22151.4 19159.2
écart-type 860.163 922.202 1095.72 1210.81 1095.53 1077.39 1326.36
class 6 (687)
moyennes 10676.7 11588.5 14467.9 17307.1 21387.2 23952.5 20855.6
écart-type 654.394 760.444 920.24 991.858 885.517 836.571 1358.26
class 7 (1084)
moyennes 11199.5 12174 15358.4 18707.2 22544.3 25563.6 22281.5
écart-type 714.358 800.778 934.065 1072.39 909.561 752.094 1730.2
class 8 (1293)
moyennes 11649.2 12696.2 16168.2 19957 23642.4 27108.3 23910.7
écart-type 633.69 705.108 785.212 794.661 693.652 613.02 1270.6
class 9 (1279)
moyennes 12307.8 13438.6 17221.3 21386.4 24872 28519.9 25276.5
écart-type 632.726 697.061 722.533 685.491 581.066 591.761 961.26
class 10 (631)
moyennes 12925.2 14153 18279.8 22826.1 26177.6 29865 26582.6
écart-type 633.357 697.001 739.746 645.763 615.265 671.366 1059.74

Distribution des classes
       1356         52        115        214        429
        687       1084       1293       1279        631
######## iteration 1 ###########
10 classes, 93.71% points stable
Distribution des classes
       1312        105        110        216        446
        700       1076       1257       1239        679
######## iteration 2 ###########
10 classes, 95.71% points stable
Distribution des classes
       1307        100        119        232        458
        714       1048       1222       1227        713
######## iteration 3 ###########
10 classes, 96.05% points stable
Distribution des classes
       1307        103        114        247        472
        724       1012       1201       1222        738
######## iteration 4 ###########
10 classes, 96.34% points stable
Distribution des classes
       1307        106        117        245        495
        749        942       1204       1215        760
######## iteration 5 ###########
10 classes, 96.78% points stable
Distribution des classes
       1307        112        118        246        519
        788        850       1190       1233        777
######## iteration 6 ###########
10 classes, 97.27% points stable
Distribution des classes
       1307        112        115        261        537
        823        772       1166       1253        794
######## iteration 7 ###########
10 classes, 97.79% points stable
Distribution des classes
       1307        112        113        269        562
        842        726       1137       1261        811
######## iteration 8 ###########
10 classes, 98.17% points stable
Distribution des classes
       1307        110        109        275        583
        855        699       1105       1266        831
########## final results #############
10 classes (convergence=98.2%)

class separability matrix

         1     2     3     4     5     6     7     8     9    10

  1      0
  2    2.3     0
  3    4.5   1.8     0
  4    5.4   1.0   1.5     0
  5    8.9   1.8   1.4   0.9     0
  6   12.0   2.4   1.5   1.7   1.0     0
  7   11.6   2.7   1.1   2.2   1.6   1.0     0
  8   15.4   3.2   1.5   2.8   2.2   1.4   0.7     0
  9   16.0   3.7   1.4   3.4   3.0   2.4   1.3   1.1     0
 10   16.2   4.1   1.6   3.9   3.6   3.2   2.1   2.0   0.9     0

class means/stddev for each band

class 1 (1307)
moyennes 7984.04 8340.86 8198.37 7487.63 7295.57 7438.61 7453.3
écart-type 440.074 481.009 735.788 580.503 432.552 225.643 190.224
class 2 (110)
moyennes 8928.21 9560.78 11310.4 12004.5 18804.7 14720.7 11994.4
écart-type 801.858 968.1 1359.59 2035.45 3487.33 2321.34 1896.19
class 3 (109)
moyennes 13067.6 14279.2 17805.4 21287.2 23357.9 21156.4 14672.2
écart-type 1562.46 1394.03 1573.66 2195.67 2583.4 4187.53 2631.53
class 4 (275)
moyennes 9891.63 10678.9 12890.9 14678.9 19079.4 19469.2 16407.4
écart-type 695.289 788.872 1044.94 1410.6 1778.39 1171.9 1416.14
class 5 (583)
moyennes 10293.2 11137 13711.8 16103.7 20394.6 22471 19542.5
écart-type 604.698 689.38 871.573 975.985 1124.77 811.76 961.218
class 6 (855)
moyennes 10612.2 11499.4 14448.2 17499.7 21692.3 24780.8 21881.5
écart-type 409.838 471.056 570.623 628.588 807.797 751.201 940.758
class 7 (699)
moyennes 11656.8 12699.9 15953.9 19240.5 22785.9 25621.8 22069.8
écart-type 545.254 594.096 695.762 862.717 848.127 837.002 1131.5
class 8 (1105)
moyennes 11285.4 12284.6 15704.5 19542.1 23411.3 27124.8 24314.1
écart-type 488.064 542.752 585.104 600.531 646.64 611.941 798.085
class 9 (1266)
moyennes 12272.3 13399.6 17130.4 21196.4 24648.7 28249.1 24935.7
écart-type 527.041 577.33 585.804 587.309 581.741 691.399 1118.4
class 10 (831)
moyennes 12909.5 14126.9 18200.5 22681.4 26019.6 29650.3 26310.5
écart-type 594.12 657.206 705.143 652.912 642.377 739.307 1140.6

#################### CLASSES ####################
10 classes, 98.17% points stable

######## CLUSTER END (Fri Jul 21 15:37:15 2023) ########
\end{verbatim}
\end{scriptsize}

\subsubsection[\appendixname~\thesubsubsection]{July 2023}
\begin{scriptsize}
\begin{verbatim}
#################### CLUSTER (Fri Jul 21 15:46:05 2023) ####################
Location: Tunisia_Gabes
Mapset:   PERMANENT
Group:    L8_2023j
Subgroup: res_30m
 L8_2023j_01@PERMANENT
 L8_2023j_02@PERMANENT
 L8_2023j_03@PERMANENT
 L8_2023j_04@PERMANENT
 L8_2023j_05@PERMANENT
 L8_2023j_06@PERMANENT
 L8_2023j_07@PERMANENT
Result signature file: cluster_L8_2023j

Region
  North:   3946245.00  East:    696915.00
  South:   3713085.00  West:    472455.00
  Res:          30.00  Res:         30.00
  Rows:          7772  Cols:         7482  Cells: 58150104
Mask: no

Cluster parameters
  Nombre de classes initiales: 	10
 Minimum class size:           17
 Minimum class separation:     0.000000
 Percent convergence:          98.000000
 Maximum number of iterations: 30
 Row sampling interval:        77
 Col sampling interval:        74

Sample size: 7140 points

means and standard deviations for 7 bands

moyennes 10900 11916.6 14769.6 17587.6 20785.4 23325.9 20657.7
écart-type 1727.35 1941.19 3033.77 4643.85 5867.69 7114.07 6057.61

initial means for each band

classe 1    9172.69 9975.43 11735.8 12943.8 14917.7 16211.9 14600.1
classe 2    9556.55 10406.8 12410 13975.8 16221.7 17792.8 15946.3
classe 3    9940.41 10838.2 13084.2 15007.7 17525.6 19373.7 17292.4
classe 4    10324.3 11269.6 13758.3 16039.7 18829.5 20954.6 18638.5
classe 5    10708.1 11700.9 14432.5 17071.7 20133.5 22535.5 19984.7
classe 6    11092 12132.3 15106.7 18103.6 21437.4 24116.4 21330.8
classe 7    11475.8 12563.7 15780.9 19135.6 22741.3 25697.3 22676.9
classe 8    11859.7 12995.1 16455 20167.6 24045.3 27278.2 24023.1
classe 9    12243.5 13426.4 17129.2 21199.5 25349.2 28859.1 25369.2
classe 10   12627.4 13857.8 17803.4 22231.5 26653.1 30440 26715.3

class means/stddev for each band

class 1 (1355)
moyennes 8075.49 8727.41 9437.09 9138.57 9359.69 9532.81 9286.77
écart-type 790.364 844.767 1011.44 1144.67 1646.73 1292.92 1008.97
class 2 (37)
moyennes 9332.89 10030 11943.7 12806.4 19781.9 17731 14461.7
écart-type 742.374 783.927 867.001 1355.42 2400.84 866.43 1250.87
class 3 (97)
moyennes 9814.24 10655.3 12919.1 14516.4 19582.8 19454 16223.7
écart-type 991.918 1063.54 1274.17 1700.9 1747.77 1418.78 1627.08
class 4 (222)
moyennes 10123.8 11025.9 13440.8 15478.3 19985.8 21366.4 18287.3
écart-type 734.52 821.29 1029.02 1196.54 1351.57 1092.62 1341.32
class 5 (440)
moyennes 10452.7 11403.9 14063.8 16552.6 20759.6 22946.8 19905.2
écart-type 833.424 901.147 1110.89 1235.44 1052.49 1154.96 1406.44
class 6 (776)
moyennes 10750.7 11746.6 14693.4 17651.7 21816.3 24587.3 21514.1
écart-type 811.417 865.831 1032.08 1114.29 924.035 1106.57 1369.82
class 7 (1147)
moyennes 11174.5 12242.5 15476.3 18869.8 22808.5 26040.1 22840.9
écart-type 654.23 754.414 882.666 926.683 807.241 712.3 1434.36
class 8 (1232)
moyennes 11708.2 12849.1 16350.5 20164.9 23941.3 27415.8 24110.8
écart-type 633.646 715.82 825.969 861.388 700.455 604.481 1346.38
class 9 (1127)
moyennes 12374.1 13568.6 17326.2 21497.8 25107.3 28668 25333.7
écart-type 600.598 662.628 721.794 732.328 588.117 558.394 1085.92
class 10 (707)
moyennes 13027.2 14299.3 18354.8 22884.4 26366 30006.2 26643.1
écart-type 560.819 625.236 691.516 651.408 607.079 672.174 1139.74

Distribution des classes
       1355         37         97        222        440
        776       1147       1232       1127        707
######## iteration 1 ###########
10 classes, 94.02% points stable
Distribution des classes
       1332         65         96        237        450
        788       1116       1223       1109        724
######## iteration 2 ###########
10 classes, 97.20% points stable
Distribution des classes
       1329         65        113        249        457
        783       1100       1207       1107        730
######## iteration 3 ###########
10 classes, 97.45% points stable
Distribution des classes
       1329         74         94        285        467
        766       1088       1196       1109        732
######## iteration 4 ###########
10 classes, 97.68% points stable
Distribution des classes
       1329         83        100        289        491
        757       1074       1172       1112        733
######## iteration 5 ###########
10 classes, 97.54% points stable
Distribution des classes
       1329         90        105        296        507
        759       1056       1154       1105        739
######## iteration 6 ###########
10 classes, 98.10% points stable
Distribution des classes
       1330         92        110        294        531
        762       1043       1131       1102        745
########## final results #############
10 classes (convergence=98.1%)

class separability matrix

         1     2     3     4     5     6     7     8     9    10

  1      0
  2    2.1     0
  3    3.3   1.8     0
  4    3.8   1.0   1.5     0
  5    5.2   1.7   1.4   0.8     0
  6    6.4   2.3   1.5   1.6   0.8     0
  7    6.7   2.6   1.2   2.1   1.4   0.8     0
  8    7.9   3.3   1.3   2.9   2.3   1.8   0.8     0
  9    8.1   3.6   1.2   3.4   2.9   2.6   1.6   1.0     0
 10    8.7   4.2   1.4   4.1   3.7   3.5   2.5   1.9   0.9     0

class means/stddev for each band

class 1 (1330)
moyennes 8058.7 8709.53 9401.52 9085.59 9207.23 9425.65 9225.7
écart-type 781.954 834.63 972.584 1066.02 1181.54 1015.2 891.046
class 2 (92)
moyennes 9291.92 10046.7 12015.2 12970.4 19432.1 17349.3 14201.8
écart-type 745.954 838.223 1129.21 1705.72 2561.57 1726.6 1609.7
class 3 (110)
moyennes 13203.6 14561.9 18278.6 21909.7 24091.4 22330.7 16271.9
écart-type 1567.76 1476.23 1785.22 2369.29 2567.26 3546.38 2580.62
class 4 (294)
moyennes 9988.86 10863.6 13228.3 15183.8 19847.5 21198.4 18095.5
écart-type 591.177 646.475 832.944 1064.72 1442.62 991.638 1231.61
class 5 (531)
moyennes 10442.8 11398 14076.8 16579.4 20872.4 23298.3 20289.1
écart-type 585.83 656.622 797.926 822.723 1031.1 727.682 878.071
class 6 (762)
moyennes 10610.1 11580.5 14528 17610 21880.5 25080.9 22128.3
écart-type 410.133 480.267 555.304 529.062 738.722 624.364 791.019
class 7 (1043)
moyennes 11216.5 12291.5 15544.7 18939.2 22856.3 26161.6 23039.6
écart-type 538.054 605.063 670.903 642.888 680.597 720.558 1047.69
class 8 (1131)
moyennes 11576.7 12699.9 16193.5 20040.4 23904.6 27593 24533.7
écart-type 489.157 546.991 572.443 528.477 548.033 588.183 875.682
class 9 (1102)
moyennes 12457.3 13664 17416.9 21541.8 25080.1 28529 25049.1
écart-type 505.335 558.335 603.626 628.159 559.537 742.932 1348.62
class 10 (745)
moyennes 13039.2 14312.6 18361.6 22877.6 26335.2 29946.2 26566.6
écart-type 546.893 609.428 674.513 637.603 613.11 704.752 1167.04

#################### CLASSES ####################

10 classes, 98.10% points stable

######## CLUSTER END (Fri Jul 21 15:46:05 2023) ########
\end{verbatim}
\end{scriptsize}

\let\cleardoublepage\clearpage

%%%%%%%%%%%%%%%%%%%%%%%%%%%%%%%%%%%%%%%%%%
\begin{adjustwidth}{-\extralength}{0cm}
%\printendnotes[custom] % Un-comment to print a list of endnotes

\reftitle{References}
%\nocite{*}

%=====================================
% References, variant A: external bibliography
%=====================================
\bibliography{Tunisia.bib}
%\nocite{*}

%%%%%%%%%%%%%%%%%%%%%%%%%%%%%%%%%%%%%%%%%%
\end{adjustwidth}
\end{document}
